\documentclass{article}

% ------------------------------------------------------------------
% NeurIPS 2026 style. Default (anonymous) option while under review.
% ------------------------------------------------------------------
\usepackage{neurips_2026}

\usepackage[utf8]{inputenc}
\usepackage[T1]{fontenc}
\usepackage{hyperref}
\usepackage{url}
\usepackage{booktabs}
\usepackage{amsmath}
\usepackage{amssymb}
\usepackage{amsthm}
\usepackage{mathtools}
\usepackage{microtype}
\usepackage{graphicx}
\usepackage{enumitem}
\usepackage{natbib}

\theoremstyle{plain}
\newtheorem{theorem}{Theorem}
\newtheorem{proposition}{Proposition}
\newtheorem{lemma}{Lemma}
\newtheorem{corollary}{Corollary}

\theoremstyle{definition}
\newtheorem{definition}{Definition}
\newtheorem{assumption}{Assumption}
\newtheorem{prediction}{Prediction}
\newtheorem{example}{Example}

\theoremstyle{remark}
\newtheorem{remark}{Remark}

\newcommand{\E}{\mathbb{E}}
\newcommand{\Prob}{\mathbb{P}}
\newcommand{\Var}{\operatorname{Var}}
\newcommand{\R}{\mathbb{R}}
\newcommand{\1}{\mathbf{1}}
\newcommand{\cD}{\mathcal{D}}
\newcommand{\cH}{\mathcal{H}}
\newcommand{\cL}{\mathcal{L}}
\newcommand{\cU}{\mathcal{U}}
\newcommand{\cF}{\mathcal{F}}
\newcommand{\indep}{\perp\!\!\!\perp}
\newcommand{\good}{\mathrm{good}}
\newcommand{\bad}{\mathrm{bad}}
\newcommand{\Tr}{\operatorname{tr}}
\newcommand{\vecop}{\operatorname{vec}}
\newcommand{\Om}{\Omega}
\newcommand{\wdc}{\mathrm{wd}}
\newcommand{\Adam}{\mathrm{AdamW}}

% A short plain-language summary after a formal statement.
\newcommand{\takeaway}[1]{%
  \par\medskip\noindent\textit{In words:} #1\par\medskip}

\title{A Tipping Point for Noisy Think Traces,\\
and Why It Moves When You Change the Batch Size}

\author{%
  Anonymous Author(s)\\
  Affiliation\\
  Address\\
  \texttt{email}
}

\begin{document}

\maketitle

% ==================================================================
\begin{abstract}
Train a small transformer on chain-of-thought ``think traces'' of which only a
fraction $\rho$ are correct, and accuracy does not improve smoothly with $\rho$.
It sits at a floor, then jumps. We take one synthetic task apart far enough to
say what the jump is.

The first thing we find is that it is not about information. In our task the
trace is conditionally independent of the answer given the prompt, so no
predictor can use it to do better. The ceiling and the floor are both available
in closed form, $0.893$ and $0.157$, and the measured curve runs between them
with nothing fitted. Whatever the transition is, it is about which computation
the optimizer lands on, not about what is learnable.

So we look at the optimizer. A single scalar parameter---how much the answer
reads off the trace---feels a drift proportional to $\rho$ and a minibatch noise
proportional to $B^{-1/2}$. The transition is where the drift wins. Because each
example flips its own correct/corrupted coin, the whole mixture variance sits
inside the $1/B$, and the critical fraction is $\rho_c(B)=\kappa\sqrt{v_+}/
(\mu_0\sqrt B)+O(B^{-1})$. There is no fixed critical fraction: it vanishes with
batch size. The drift constant is not a free parameter either, since
$\mu_0\propto 0.893-0.157$ and the leftover scale cancels.

A three-seed sweep gets the shape right and the rate wrong. The threshold falls
by $7.7\times$ over a $16\times$ range of $B$, but the measured exponent is
$-0.72\pm0.05$ where we predicted $-1/2$. We chase the two obvious culprits.
AdamW changes the mechanism in several interesting ways and leaves the exponent
at $-1/2$; parameter dimension, in a linear-attention surrogate, does the same.
What survives is total sample exposure, which we quantify with a susceptibility
identity and isolate with a one-variable control. We also report an anomaly at
$\rho=1$ that appears in all fifteen runs, and trace a likely cause to a token
budget in our own evaluation script.
\end{abstract}

% ==================================================================
\section{The phenomenon, and how to read this paper}
\label{sec:intro}

Chain-of-thought training puts a loss on intermediate reasoning tokens as well as
the final answer. The traces used for that supervision are rarely all good. Some
are valid derivations, some wander, and in a distilled pipeline a real fraction
are simply wrong. This raises a concrete question: as the fraction of good traces
goes up, when does the model start \emph{using} the trace as scaffolding instead
of ignoring it?

The answer, at least in small settings, is not ``gradually''. Here is one row of
our data, at batch size $128$:

\begin{center}
\begin{tabular}{lccccc}
\toprule
$\rho$ (fraction of correct traces) & 0.10 & 0.20 & 0.30 & 0.40 & 0.50\\
accuracy (\%) & 17.8 & 27.2 & 47.9 & 81.6 & 87.9\\
\bottomrule
\end{tabular}
\end{center}

Between $\rho=0.2$ and $\rho=0.4$ the model goes from barely above guessing to
near its ceiling. That is the object of study. What is it?

\paragraph{Three suspects.}
Only three kinds of explanation are available, and it is worth naming them before
any math.

\begin{enumerate}[leftmargin=*,itemsep=3pt]
\item \textbf{Information.} Maybe correct traces tell the model something the
prompt does not. Then low $\rho$ is a data problem, and the transition is the
point where enough signal is present to learn from.
\item \textbf{Optimization.} Maybe the information is the same either way, and
what changes is whether gradient descent finds a particular circuit before noise
scrambles the attempt. Then the transition is a signal-to-noise event, and it
should move when we change anything that changes the noise.
\item \textbf{Representability.} Maybe the trace-using circuit does not fit in the
architecture at low $\rho$. This is the least likely of the three, and we deal
with it briefly in Appendix~\ref{app:repr}.
\end{enumerate}

We rule out the first, develop the second into a prediction we can test, and set
the third aside.

\paragraph{Suspect 1 is innocent, and that is the useful part.}
Our task asks the model to find which occurrence of a character generated a
query. Section~\ref{sec:bayes} shows the trace is conditionally independent of
the answer given the prompt. The argument is one sentence long once you see it:
the correct trace is a deterministic function of the word, so it cannot tell you
anything about the word you did not already have, and the corrupted trace is
built out of fresh randomness that has nothing to do with the answer.

This costs us the most tempting story and buys us something better. Because the
task is fully solvable on paper, we know exactly where the ceiling is
($89.3\%$, set by how often a character repeats in a random word) and exactly
where the floor is ($15.7\%$, guessing among positions). Every accuracy in this
paper can be read as a fraction of the distance between two computed numbers,
with nothing fitted. Papers about sharp transitions usually cannot do that.

\paragraph{Suspect 2, made quantitative.}
Section~\ref{sec:scalar} strips the model down to one number: a scalar $u$
controlling how much the answer logits listen to the trace. Correct traces push
$u$ up at a rate proportional to $\rho$. All traces, correct and corrupted, add
minibatch noise of size $B^{-1/2}$. The transition is where push beats jitter:
\begin{equation}
\label{eq:intro-threshold}
\rho_c(B,\kappa) \;=\; \frac{\kappa\sqrt{v_+}}{\mu_0\sqrt B} + O(B^{-1}).
\end{equation}
Read that as a claim with teeth. It says there is no such thing as ``the''
critical trace fraction. Quadruple the batch and the threshold halves. It also
says the constants are not ours to tune: $\mu_0\propto 0.893-0.157$ comes
straight from the combinatorics of the task, and the one genuinely unknown scale
cancels between numerator and denominator (Remark~\ref{rem:scale}).

\paragraph{What happened when we ran it.}
The shape survives: over three seeds, $\rho_c$ drops by $7.7\times$ as $B$ grows
$16\times$, monotonically, in every seed. The rate does not: we measure
$-0.72\pm0.05$, not $-1/2$.

Most of the second half of the paper is about that gap. Two objections to our
derivation are obvious---it is written for plain SGD but we ran AdamW, and it is
written for one coordinate but the model has $8.1\times10^5$ of them---so we work
both through. AdamW turns out to change the story in ways worth knowing about
(the drift saturates; the per-step noise stops shrinking with $B$ and the drift
grows instead; weight decay converts a ``did it escape in time'' criterion into a
horizon-free one), and yet the exponent comes out at $-1/2$ again. Dimension, in
a linear-attention surrogate, does not touch the exponent either. Two further
mechanisms we can name push the exponent the wrong way, toward $0$.

What is left standing is total sample exposure: at a fixed number of steps, a
bigger batch also means more data seen. We give the identity that converts a
measured exponent into a required sample-scaling exponent
(Proposition~\ref{prop:susceptibility}: $0.22$ would do it) and design the
one-variable experiment that would test it (Prediction~\ref{pred:diversity}).
We also leave open the dull possibility that three seeds is too few, since a
seed-level confidence interval on the exponent does contain $-1/2$.

\paragraph{How to read the rest.}
Sections~\ref{sec:setup} and~\ref{sec:bayes} are self-contained and elementary;
if you only want to know what the task is and why the trace is useless, stop
after them. Section~\ref{sec:scalar} is the core argument and the only place the
main prediction is derived. Sections~\ref{sec:adam} and~\ref{sec:matrix} are
defensive: they exist to answer ``but you used Adam'' and ``but your model is not
one-dimensional'', and can be skipped on a first pass.
Sections~\ref{sec:experiments}--\ref{sec:gap} are the measurements and the
autopsy. Throughout we mark plainly which statements are theorems from stated
hypotheses, which are surrogate approximations, which are borrowed, and which are
things we measured. Table~\ref{tab:notation} collects the symbols, and is worth a
glance before Section~\ref{sec:scalar}.

\subsection{Contributions}

\begin{enumerate}[leftmargin=*,itemsep=2pt]
\item A synthetic think-trace task specified tightly enough that the Bayes
ceiling and the chance floor are both closed forms, plus a proof that the trace
is Bayes-irrelevant in it (Section~\ref{sec:bayes}).
\item A finite-batch escape-and-persistence theorem for a scalar trace-copy
coordinate, with the mixture variance accounted for correctly and a Freedman
bound whose deviation scale is matched to the drift slack. This yields
$\rho_c=\Theta(B^{-1/2})$, with the drift constant derived from the task and the
coordinate scale shown to cancel (Section~\ref{sec:scalar}).
\item A quasi-stationary analysis of AdamW on the same coordinate, giving a
saturating drift field and the same leading exponent, plus a fixed-horizon
invariance theorem covering scale-invariant optimizers generally
(Section~\ref{sec:adam}).
\item A lift of the argument to the parameter manifold of a linear attention
layer, showing where dimension enters and why it misses the threshold
(Section~\ref{sec:matrix}).
\item A three-seed batch-size sweep, a measured exponent of $-0.72\pm0.05$, a
susceptibility identity relating the residual to sample exposure, and a
token-budget confound in our own evaluation protocol that bears on the $\rho=1$
anomaly (Sections~\ref{sec:experiments} and~\ref{sec:gap}).
\end{enumerate}

\begin{table}[t]
\caption{Notation used throughout. Symbols local to one section are defined
there.}
\label{tab:notation}
\centering
\small
\begin{tabular}{ll}
\toprule
Symbol & Meaning \\
\midrule
$\Sigma, q$ & alphabet and its size ($q = 26$) \\
$L, \ell$ & word length (random / fixed) \\
$W, C, I^\star$ & word, queried character, gold index; the target is $Y = I^\star$ \\
$K(W,C)$ & multiplicity: how many $j$ have $W_j = C$ \\
$R$, $Z$ & the presented trace; indicator that it is the correct one \\
$\rho$, $B$, $T$ & correct-trace fraction, batch size, number of steps \\
$u$, $c$ & scalar trace-copy coordinate; its gradient scale \eqref{eq:gradform} \\
$\mu_0,\mu_1$ & lower / upper bound on the drift from correct traces \\
$\epsilon_\bad,\bar\epsilon$ & bias from corrupted traces, and its bound \\
$v_+$, $V$ & per-example gradient variance bound; mixture bound $v_+ +\rho(1-\rho)\Delta^2$ \\
$\kappa$, $r$ & confidence factor; standardized SNR $\rho\mu_0\sqrt{B/V}$ \\
$\bar A^\star$, $A_{\mathrm{chance}}$ & Bayes accuracy $\approx0.893$; chance $\E[1/L]\approx0.157$ \\
$S_T$ & cumulative step size $\sum_{s\le T}\eta_s$ \\
\bottomrule
\end{tabular}
\end{table}

% ==================================================================
\section{Related work}
\label{sec:related}

\paragraph{Induction heads and copying circuits.}
The circuit we care about is a close relative of the induction head, which
completes patterns $A\,B\ldots A\to B$ \citep{olsson2022incontext}.
\citet{nichani2024how} analyse how a two-layer transformer picks up causal
structure by gradient descent, including induction-head-like dynamics; our scalar
argument is in that spirit, specialised to a mixture of correct and corrupted
supervision. \citet{chen2024unveiling} give provable training dynamics with an
explicit feature-learning phase, and \citet{ekbote2025what} prove a depth
separation between one and two layers for induction heads on Markov chains, which
we adopt rather than re-derive (Appendix~\ref{app:repr}).

\paragraph{Phase transitions and emergence.}
\citet{lavie2026phase} develop a Bayesian theory of copy-head emergence with
phase transitions driven by task statistics. The debate over emergent abilities
\citep{wei2022emergent,schaeffer2023mirage} is why we insist on pinning down both
reference levels before calling anything sharp: with a computed floor and a
computed ceiling we can say what fraction of the available range a transition
crosses, instead of reporting a raw jump whose size depends on the metric.

\paragraph{Sample complexity and information exponents.}
That the required sample size blows up as the drift vanishes is familiar from
work on information exponents for single- and multi-index models
\citep{benarous2021online,damian2023smoothing,abbe2023sgd}, where escaping a
symmetric initialization needs samples scaling as a negative power of the
population correlation. Appendix~\ref{app:samplecomplexity} states the version
that holds here, and shows it is the same inequality as our threshold seen from
another angle.

\paragraph{Learning under label noise.}
We are learning from a mixture of clean and corrupted supervision, but not in the
classical label-noise sense \citep{natarajan2013learning}: the corruption leaves
the final answer token alone and damages only auxiliary trace tokens. Since the
corrupted trace is independent of the answer by construction, it contributes zero
\emph{initial} bias to the gradient statistic we track. That is exactly what makes
the threshold depend on batch size rather than on the noise rate.

\paragraph{Adaptive optimizers and batch-size scaling.}
Section~\ref{sec:adam} works with AdamW
\citep{kingma2015adam,loshchilov2019decoupled}. Treating the accumulators as
quasi-stationary follows the SDE approximations and square-root scaling rules of
\citet{malladi2022sdes}, and the saturating drift we obtain interpolates between
gradient descent and sign descent \citep{bernstein2018signsgd};
\citet{zhang2020adaptive} give a complementary account of why adaptivity helps
attention models. The closest existing framing is the critical batch size beyond
which gradient noise stops binding \citep{mccandlish2018empirical}; our
standardized SNR is that idea localized to a single circuit under mixture
supervision. The sample-exposure channel we cannot rule out lives in the regime
described by data-scaling laws \citep{kaplan2020scaling}.

\paragraph{Expressivity of chain-of-thought, and linear attention.}
A separate literature asks what intermediate tokens let a fixed-depth transformer
compute \citep{feng2023towards,merrill2024expressive,li2024chain}. Our task is
built so that this question cannot arise: by Proposition~\ref{prop:indep} the
trace carries no information about the answer beyond the prompt.
Section~\ref{sec:matrix} uses the softmax-free attention map
\citep{katharopoulos2020transformers}, the usual vehicle for provable statements
about attention dynamics \citep{vonoswald2023transformers,zhang2024incontext},
with dimensional factors from \citet{tropp2011freedman}.

% ==================================================================
\section{The task, starting from one example}
\label{sec:setup}

Before any definitions, here is the whole task in one instance.

\begin{example}[The task in one line]
\label{ex:instance}
The word is \texttt{banana}. We pick a position, say $I^\star=3$, and read off
the character sitting there: $C=\texttt{a}$. The model sees the prompt
\texttt{banana;a} and must answer $3$.

Notice immediately what the model is up against. The letter \texttt{a} occurs at
positions $1$, $3$ and $5$, and nothing in the prompt distinguishes them. The
best any predictor can do on this instance is $1/3$.

Now the trace. The correct trace is \texttt{0b 1a 2n 3a 4n 5a}: a positionally
annotated copy of the word. Does it help? No. It shows the same three
\texttt{a}'s, with no marker on the one we chose. It is a re-encoding of
information the model already had.
\end{example}

That last observation is the whole of Section~\ref{sec:bayes}, and it is worth
sitting with, because it is a design choice rather than an accident. We could
have made the target $\min\{j:W_j=C\}$, in which case the answer would be a
deterministic function of the prompt and Bayes accuracy would be $100\%$. Instead
we draw the position first and read the query off it, which leaves a genuine
$1/K$ ambiguity and gives us a ceiling below $1$ to measure against.

\begin{definition}[Word-index task]
\label{def:task}
Draw $L$; draw $W = (W_0,\dots,W_{L-1})$ uniformly from $\Sigma^L$; draw
$I^\star\sim\mathrm{Unif}\{0,\dots,L-1\}$ independently of $W$ given $L$; set
$C := W_{I^\star}$ and the prompt $P := (W;C)$; draw a trace $R$ per
Definition~\ref{def:traces}. The target is $Y := I^\star$.
\end{definition}

In experiments $\Sigma$ is the lowercase Latin alphabet, so $q=26$, and
$L\sim\mathrm{Unif}\{4,\dots,10\}$.

\begin{definition}[Correct and corrupted traces]
\label{def:traces}
The correct trace is $R^{\good}(W) = (0,W_0)(1,W_1)\cdots(L-1,W_{L-1})$. The
corrupted trace is
$R^{\bad} = (\Pi(0),\widetilde W_0)\cdots(\Pi(L-1),\widetilde W_{L-1})$, where
$\Pi$ is a uniform random permutation and the $\widetilde W_j$ are i.i.d.\
uniform on $\Sigma$, all independent of $W$, of $I^\star$, and of each other.
\end{definition}

The corruption is deliberately \emph{shape-preserving}. Same length, same number
of segments, same alternation of index and character fields, same marginal
distribution of index values. The only thing missing is any relationship between
the two fields. A model cannot spot a corrupted trace from surface statistics; it
would have to check the trace against the word.

\begin{definition}[Training mixture]
\label{def:mixture}
For $\rho\in[0,1]$, draw $Z\sim\mathrm{Bernoulli}(\rho)$ \emph{independently for
each example} and set $R = R^{\good}$ if $Z=1$ and $R=R^{\bad}$ otherwise. Write
$\cD_\rho$ for the resulting joint law.
\end{definition}

The italics matter more than they look. Section~\ref{sec:variance} shows that
per-example independence of $Z$ is precisely what puts the between-component
variance inside the $1/B$, and hence what makes the threshold vanish with batch
size. It is also easy to break by accident, by building homogeneous batches or
shuffling at shard level, so Section~\ref{sec:limitations} lists it as a scope
condition rather than a technicality.

\paragraph{Sequence format and loss.}
A training sequence is $W\texttt{;}C\ R\ \texttt{:}\ I^\star\
\texttt{\textbackslash n}$. With prompt masking, the loss covers every token
after the prompt, so writing $R$ for trace positions and $Y$ for answer
positions,
\begin{equation}
\label{eq:lfull}
\cL_{\mathrm{full}}
= \frac{1}{|R| + |Y|}\Big(
\sum_{i \in R}\cL_{\mathrm{CE}}(r_i, r_i^\star)
+ \sum_{j \in Y}\cL_{\mathrm{CE}}(y_j, y_j^\star)\Big).
\end{equation}
Since $|R|\approx3\ell$ and $|Y|\approx1$, roughly $95\%$ of the supervised
tokens are trace tokens. Section~\ref{sec:competing} takes that seriously as a
rival explanation for everything in this paper.

\paragraph{Two ways to solve the task.}
It helps to have names for the two algorithm families. Write
$\cH_{\mathrm{short}}$ for predictors that search the prompt directly for
positions $j$ with $W_j=C$ and then break ties somehow, and $\cH_{\mathrm{trace}}$
for predictors that find a trace segment whose character field matches $C$ and
copy that segment's index into the answer.

The point of the next section is that these two families are \emph{tied}. They
have the same population ceiling. So the question is never which is better; it is
which one a finite training run with a finite batch size happens to find, and how
that depends on $\rho$.

% ==================================================================
\section{The trace cannot help, and we can compute exactly how well anything can do}
\label{sec:bayes}

\begin{proposition}[Conditional independence]
\label{prop:indep}
Under Definitions~\ref{def:task}--\ref{def:mixture}, for every $\rho$,
$I^\star \indep R \mid (W,C)$.
\end{proposition}

\begin{proof}
\emph{Step 1: $I^\star\indep R\mid W$.} Condition on $W$ and $L$. The correct
trace is a deterministic function of $W$, so it is conditionally independent of
everything given $W$. The corrupted trace is a function of $(\Pi,\widetilde W)$
alone---notation notwithstanding, it does not depend on $W$---and
$(\Pi,\widetilde W)\indep I^\star\mid W$. The indicator $Z$ is independent of
everything else. So $R = \varphi(W,\Pi,\widetilde W,Z)$ where every argument
other than $W$ is conditionally independent of $I^\star$ given $W$, which gives
$I^\star\indep R\mid W$.

\emph{Step 2: adding $C$.} This is not automatic, because $C=W_{I^\star}$ is a
function of $(W,I^\star)$, so conditioning on $C$ genuinely changes the
conditional law of $I^\star$. Step 1 gives
$\E[f(R)g(I^\star)\mid W] = \E[f(R)\mid W]\,\E[g(I^\star)\mid W]$ for bounded
$f,g$. Conditioning further on $C$ restricts $I^\star$ to $\{j:W_j=c\}$ and
reweights its law. But the conditional law of $R$ given $W$ never depended on
$I^\star$ in the first place, so
$\Prob(R\in\cdot\mid W,C=c) = \Prob(R\in\cdot\mid W)$ for every $c$, and the
factorisation survives.
\end{proof}

\begin{corollary}[Trace irrelevance]
\label{cor:irrelevance}
$\Prob(I^\star = j\mid W,C,R) = \Prob(I^\star = j\mid W,C)$ for every $j$ and
every $\rho$. The Bayes-optimal predictor from $(W,C,R)$ is the Bayes-optimal
predictor from $(W,C)$, with identical $0$--$1$ risk.
\end{corollary}

\takeaway{Handing the model a perfect trace is worth exactly nothing at the
population level. This is why you will not find a theorem here of the form ``the
trace-using predictor has strictly lower risk''. In this task that statement is
false.}

\paragraph{Where the ceiling comes from.}
Fix $W$ and $C$ and let $K=|\{j:W_j=C\}|\ge1$. By construction
$\Prob(I^\star=j\mid W,C) = \1\{W_j=C\}/K$: the gold index is uniform over the
matching positions. So the Bayes predictor gets $1/K$, and it gets that by any
rule that commits to a single matching position. ``Pick the first match'' is
Bayes-optimal, which is how the prompt-only family $\cH_{\mathrm{short}}$ can be
optimal despite doing nothing clever.

The ceiling is therefore a question about random words: how often does a
character repeat?

\begin{proposition}[Closed-form Bayes accuracy]
\label{prop:bayesacc}
For a word of length $\ell$ drawn uniformly from $\Sigma^\ell$,
\begin{equation}
\label{eq:bayesacc}
A^\star(\ell) \;=\; \E[1/K] \;=\; \frac{q}{\ell}\Big[1 - \big(1-\tfrac1q\big)^{\ell}\Big].
\end{equation}
\end{proposition}

\begin{proof}
Condition on $I^\star = i$. The other $\ell-1$ characters are i.i.d.\ uniform and
independent of $C=W_i$, so the number of \emph{other} occurrences is
$X\sim\mathrm{Binomial}(\ell-1,p)$ with $p=1/q$, and $K=1+X$. The identity
$\frac{1}{1+x}\binom{\ell-1}{x} = \frac1\ell\binom{\ell}{x+1}$ (write both sides
in factorials) turns the awkward expectation into a binomial sum:
\[
\E\Big[\frac1{1+X}\Big]
= \frac1\ell\sum_{x=0}^{\ell-1}\binom{\ell}{x+1}p^{x}(1-p)^{\ell-1-x}
= \frac{1}{\ell p}\Big[\sum_{k=0}^{\ell}\binom{\ell}{k}p^{k}(1-p)^{\ell-k} - (1-p)^\ell\Big]
= \frac{1-(1-p)^{\ell}}{\ell p}.
\qedhere
\]
\end{proof}

\begin{corollary}[The two numbers we measure against]
\label{cor:levels}
For $q = 26$ and $L\sim\mathrm{Unif}\{4,\dots,10\}$,
\begin{equation}
\label{eq:levels}
\bar A^\star = \tfrac17\textstyle\sum_{\ell=4}^{10} A^\star(\ell) \approx 0.893,
\qquad
A_{\mathrm{chance}} = \E[1/L] \approx 0.157 .
\end{equation}
\end{corollary}

Per-length, $A^\star(\ell)$ runs $0.944, 0.926, 0.909, 0.892, 0.875, 0.859,
0.844$ for $\ell=4,\dots,10$: longer words repeat characters more often and are
harder.

The floor deserves a remark of its own. $A_{\mathrm{chance}}=\E[1/L]$ is what you
get by guessing uniformly among the $L$ positions, which is what a model that has
learned only the output format achieves. It also happens to equal the accuracy of
the best \emph{constant} predictor, because $\Prob(I^\star=0)=\E[1/L]$ and
position $0$ is legal at every length. The two coincide here and come apart in
other tasks (Appendix~\ref{app:family}), where the constant-predictor version is
the right one to use---a model that ignores its input is the thing you want to
beat.

\begin{remark}[If you want the trace to actually matter, change the task]
\label{rem:gap}
Everything above follows from $R^{\good}$ being a function of $W$ alone. To build
a task where trace access genuinely raises the ceiling, the correct trace has to
depend on $I^\star$---it has to contain the answer, the way a real derivation
does---while corruption destroys that dependence. Shrinking the alphabet makes
the gap dramatic: at $q=4$, $\ell=10$ the prompt-only ceiling is about $0.377$,
leaving $62$ points for a trace to recover. We think this is the single most
valuable follow-up to this work.
\end{remark}

% ==================================================================
\section{One knob: a scalar theory of the threshold}
\label{sec:scalar}

We now build the argument that produces \eqref{eq:intro-threshold}. The plan is
to reduce the model to a single number, watch that number get pushed around by
two competing forces, and ask when the push wins.

\subsection{The knob, and a sign convention}

Let $u$ be a scalar parameter that enters the model only through a logit at the
answer position, in such a way that raising $u$ raises the log-odds of the index
encoded in the attended trace segment. Think of it as a volume dial on
trace-copying.

Training is $u_{t+1} = u_t - \eta_t\bar g_t$, with
$\bar g_t = B^{-1}\sum_b \hat g_t^{(b)}$ a fresh minibatch from $\cD_\rho$ and
each $Z$ drawn independently. Because the update subtracts the gradient, all
hypotheses below are stated on $h_t := -\bar g_t$. Positive drift then means
``trace-copying is being reinforced'', which is the reading we want.

\begin{assumption}[Drift, variance and tails]
\label{ass:main}
There are a neighbourhood $\cU = (-a,a)$ of $u_0 = 0$, constants
$0<\mu_0\le\mu_1<\infty$ and $v_+ > 0$, and a function
$\epsilon_\bad:\cU\to[0,\infty)$ with $\sup_\cU\epsilon_\bad\le\bar\epsilon$,
such that for all $u\in\cU$:
\begin{enumerate}[leftmargin=*,itemsep=1pt,label=(\arabic*)]
\item correct traces push, and not too hard: $\mu_0\le\E[-\hat g^{\good}(u)]\le\mu_1$;
\item corrupted traces push only a little: $|\E[-\hat g^{\bad}(u)]|\le\epsilon_\bad(u)$;
\item both have bounded variance: $\Var(\hat g^{\good}),\Var(\hat g^{\bad})\le v_+$;
\item tails are sub-exponential:
$\E\exp(\lambda(\hat g-\E\hat g))\le\exp(\lambda^2v_+/2)$ for $|\lambda|\le1/K_e$.
\end{enumerate}
We call $\bar\epsilon = 0$ the \emph{zero-bias regime}.
\end{assumption}

Two design notes. Item (1) is two-sided on purpose: a lower bound alone proves
escape but can never prevent it, so it cannot support the persistence half of
Theorem~\ref{thm:escape}. Item (4) is stated once as a primitive so that the
escape theorem and the sample-complexity statement in
Appendix~\ref{app:samplecomplexity} can share it.

These constants would be empty if we left them free. Sections~\ref{sec:drift}
and~\ref{sec:bias} compute $\mu_0$, $v_+$ and $\epsilon_\bad$ from the task.

\subsection{Where the \texorpdfstring{$1/B$}{1/B} goes}
\label{sec:variance}

This subsection is three lines of algebra and carries the paper's main structural
claim, so it is worth slowing down for.

Each example's gradient is $\hat g = Z\hat g^{\good} + (1-Z)\hat g^{\bad}$. The
law of total variance splits its variance into two pieces:
\begin{equation}
\label{eq:totalvar}
\Var(\hat g\mid u)
= \underbrace{\rho\Var(\hat g^{\good}) + (1-\rho)\Var(\hat g^{\bad})}_{\text{within components}}
+ \underbrace{\rho(1-\rho)\big(\E\hat g^{\good}-\E\hat g^{\bad}\big)^2}_{\text{between components}} .
\end{equation}
The first term is the ordinary noise inside each kind of example. The second is
the extra scatter from not knowing which kind you drew. Writing
$\Delta := \mu_1 + \bar\epsilon$ for an upper bound on the gap between the two
means---note it is the \emph{upper} bound $\mu_1$ that belongs here, since we are
bounding a variance from above---Assumption~\ref{ass:main} gives
$\Var(\hat g\mid u)\le v_+ + \rho(1-\rho)\Delta^2 =: V$.

Now the step everything turns on. Definition~\ref{def:mixture} draws $Z$
independently for each example, so the $B$ summands in $\bar g_t$ are i.i.d.\
draws from the mixture law, and
\begin{equation}
\label{eq:batchvar}
\Var(\bar g_t\mid u_t) \;\le\; \frac{V}{B}
\;=\;\frac1B\Big[v_+ + \rho(1-\rho)\Delta^2\Big],
\end{equation}
with the between-component term \emph{inside} the $1/B$.

\paragraph{The counterfactual, so the crux is visible.}
Suppose instead that a single coin decided the whole minibatch. Then the $B$
summands would be i.i.d.\ only \emph{conditionally} on $Z$, and applying the law
of total variance to $\bar g_t$ itself would give
\[
\Var(\bar g_t) = \tfrac1B[\text{within}] + \rho(1-\rho)\Delta^2,
\]
with no $1/B$ on the second term. Averaging cannot help with a coin you only flip
once. Requiring drift to beat noise would then leave a strictly positive
threshold $\kappa^2/(1+\kappa^2)$ that no amount of batching removes.

So the qualitative prediction of this paper---threshold vanishes---is downstream
of one modelling choice about how the mixture is assembled. That is a feature: it
makes the theory refutable by an experiment on data ordering, not just on batch
size.

\subsection{Drift versus noise, and the threshold}

In the zero-bias regime the drift is $\E[h_t\mid u_t]\ge\rho\mu_0$: only the
$\rho$ fraction of correct traces contributes. Ask that it exceed the noise
standard deviation by a safety factor $\kappa$,
\[
\rho\mu_0 \;\ge\; \kappa\sqrt{V/B},
\]
square both sides and collect powers of $\rho$:
\begin{equation}
\label{eq:quadratic}
\rho^2\Big(\mu_0^2 + \tfrac{\kappa^2\mu_1^2}{B}\Big)
\;-\; \rho\,\tfrac{\kappa^2\mu_1^2}{B}
\;-\; \tfrac{\kappa^2 v_+}{B} \;\ge\; 0 .
\end{equation}
Look at the coefficient of $\rho^2$. It tends to $\mu_0^2$, a constant, because
the mixture term arrives at order $1/B$ rather than order $1$. That is what makes
the root shrink. Define $\rho_c(B,\kappa)$ as the unique nonnegative root of
\eqref{eq:quadratic} at equality; it exists and is unique because with
$A:=\mu_0^2+\kappa^2\mu_1^2/B$, $b:=\kappa^2\mu_1^2/B$ and $c_q:=\kappa^2v_+/B$
the product of the roots, $-c_q/A$, is negative.

\begin{theorem}[Finite-batch threshold scaling]
\label{thm:scaling}
If $\mu_0 > 0$, $v_+ > 0$ and $\mu_1 <\infty$, then
\begin{equation}
\label{eq:scaling}
\rho_c(B,\kappa)
= \frac{\kappa\sqrt{v_+}}{\mu_0\sqrt B}
+ \frac{\kappa^2\mu_1^2}{2\mu_0^2\,B}
+ O(B^{-3/2}),
\end{equation}
so $\rho_c = \Theta(B^{-1/2})$ and $\rho_c\to0$ as $B\to\infty$.
\end{theorem}

\begin{proof}
The nonnegative root is $\rho_c = (b+\sqrt{b^2+4Ac_q})/(2A)$. The discriminant is
$b^2+4Ac_q = 4\mu_0^2\kappa^2v_+B^{-1}(1+O(B^{-1}))$, since the $B^{-2}$ pieces
are lower order, so $\sqrt{b^2+4Ac_q} = 2\mu_0\kappa\sqrt{v_+/B}+O(B^{-3/2})$.
Expanding $1/(2A) = (2\mu_0^2)^{-1}(1-\kappa^2\mu_1^2\mu_0^{-2}B^{-1}+O(B^{-2}))$
and multiplying gives the leading term $\kappa\sqrt{v_+}/(\mu_0\sqrt B)$, the
next term $\kappa^2\mu_1^2/(2\mu_0^2B)$, and cross terms at $O(B^{-3/2})$.
\end{proof}

\takeaway{In the zero-bias regime there is no batch-size-independent critical
trace fraction. Quadruple $B$ and the threshold halves. This scaling law, not any
single number, is what we are putting at risk in Section~\ref{sec:experiments}.}

\subsection{From a per-step criterion to a statement about trajectories}
\label{sec:escape}

Drift dominance is a statement about one step. It does not say the parameter
actually leaves the neighbourhood, which additionally needs enough total step
size. The next theorem supplies both directions: a condition under which the
trajectory escapes, and a condition under which it provably does not.

\begin{theorem}[Escape and persistence]
\label{thm:escape}
Let Assumption~\ref{ass:main} hold with $\bar\epsilon = 0$, let
$S_T = \sum_{s\le T}\eta_s$ and $\sum_{s\le T}\eta_s^2\le C_\eta$, and set
$\rho^{+}_{\mathrm{esc}} := a/(\mu_0S_T)$ and
$\rho^{-}_{\mathrm{per}} := a/(\mu_1S_T)$. Starting from $u_0 = 0$:
\begin{enumerate}[leftmargin=*,itemsep=2pt]
\item[(i)] If $\rho > \rho^{+}_{\mathrm{esc}}$ then, with
$\lambda := \rho\mu_0S_T - a > 0$, the trajectory exits $\cU$ in the
trace-reinforcing direction by time $T$ with probability at least
$1 - \exp\big(-\lambda^2/[2(C_\eta V/B + K_e\lambda/3)]\big)$.
\item[(ii)] If $\rho < \rho^{-}_{\mathrm{per}}$ then, with
$\lambda' := a - \rho\mu_1S_T > 0$, the trajectory stays in $\cU$ for all
$t\le T$ with probability at least
$1 - \exp\big(-\lambda'^2/[2(C_\eta V/B + K_e\lambda'/3)]\big)$.
\end{enumerate}
\end{theorem}

The proof is in Appendix~\ref{app:escape}, and two things in it are worth
flagging here.

First, the deviation scale $\lambda$ fed to Freedman's inequality is set equal to
the actual slack between cumulative drift and escape radius. This is not
cosmetic. Plug in an unrelated scale such as $\kappa\sqrt{C_\eta V/B}$ and the
exponent comes out constant in both $B$ and $\rho$, which is to say the bound
then tells you nothing about the transition.

Second, the interval $(\rho^{-}_{\mathrm{per}},\rho^{+}_{\mathrm{esc}})$ is not a
predicted transition region. It is an \emph{undecided} region, where neither
sufficient condition applies. It is nonempty exactly when $\mu_0<\mu_1$.

\begin{remark}[Two thresholds, and when they line up]
\label{rem:twothresholds}
Theorems~\ref{thm:scaling} and~\ref{thm:escape} answer different questions and it
is easy to conflate them. $\rho_c(B,\kappa)=\Theta(B^{-1/2})$ is a per-step
signal-to-noise criterion: it depends on $B$ and not on how long you train.
$\rho^{+}_{\mathrm{esc}} = a/(\mu_0S_T)$ is a budget criterion: it depends on how
long you train and not on $B$, and says that with enough cumulative step size any
positive $\rho$ escapes eventually. The two scale together only when
$S_T\asymp a\sqrt B/(\kappa\sqrt{v_+})$. This is why our sweeps hold the number
of \emph{steps} fixed rather than epochs, and it is a confound we return to in
Section~\ref{sec:gap}.
\end{remark}

\subsection{The drift constant is not a free parameter}
\label{sec:drift}

A threshold written in unnamed constants cannot be tested, so we compute them
from the task.

Picture the readout at the answer position attending to trace segments whose
character field matches $C$, and let $J$ be the index field of whichever segment
it lands on. For a logit parameter $u$ that boosts the attended index, the
cross-entropy gradient is the familiar residual,
\begin{equation}
\label{eq:gradform}
-\hat g(u) \;=\; c\,\big[\1\{J = I^\star\} - p_u(J)\big],
\qquad c > 0,
\end{equation}
with $c=\partial\ell_J/\partial u$ the sensitivity of the attended logit to $u$.
In words: the trace-copy direction gets reinforced exactly to the extent that
copying the attended index would have been right more often than the model
currently believes it will be.

To turn that into a number we need to say how attention picks $J$. This is an
assumption, not a consequence of Definition~\ref{def:traces}.

\begin{assumption}[Segment selection]
\label{ass:selection}
On correct traces, conditional on $(W,C,R)$ the selected segment $J$ is
measurable with respect to model and input randomness independent of $I^\star$.
On corrupted traces, conditional on $(W,C)$ the selected index field $J$ is
uniform on $\{0,\dots,L-1\}$ and independent of $I^\star$.
\end{assumption}

The first clause says the readout cannot break the $K$-fold tie using information
a Bayes predictor would not have. That is Corollary~\ref{cor:irrelevance}
transported onto the selection mechanism, and it is hard to object to. The second
clause is a real restriction: a trained model could in principle select segments
using the corrupted characters or positional structure, and we are assuming it
does not. Both clauses are most believable near initialization and least
believable late in training, which bounds where the following is meant to apply.

\begin{proposition}[The initial drift is the Bayes gap]
\label{prop:mu0}
Under Assumption~\ref{ass:selection}, the segments matching $C$ in a correct
trace are exactly the $K(W,C)$ positions with $W_j=C$, so
$\Prob(J = I^\star\mid W,C) = 1/K$ and
\begin{equation}
\label{eq:gooddrift}
\E[-\hat g^{\good}(u)] = c\big(\bar A^\star - \E[p_u(J)]\big),
\qquad
\E[-\hat g^{\bad}(u)] = c\big(\E[1/L] - \E[p_u(J)]\big).
\end{equation}
At a symmetric initialization $u_0$ where $p_{u_0}(j) = 1/L$ for every $j$, the
corrupted-trace drift vanishes and
\begin{equation}
\label{eq:mu0}
\mu_0 \;=\; c\,\big(\bar A^\star - \E[1/L]\big) \;\approx\; 0.736\,c .
\end{equation}
\end{proposition}

\begin{proof}
Immediate from \eqref{eq:gradform}, Assumption~\ref{ass:selection} and
Proposition~\ref{prop:bayesacc}, substituting $\E[p_{u_0}(J)] = \E[1/L]$.
\end{proof}

\takeaway{The force pulling the model toward trace-copying is proportional to the
distance between the two horizontal lines in our plots. The ceiling and the floor
are not just measurement scaffolding; their difference is the drift.}

There is a second reading of \eqref{eq:gooddrift} worth having. The drift on
correct traces is \emph{capped by Bayes accuracy}. A correct trace does not hand
the model the answer; it hands the model the same $1/K$ ambiguity the prompt
already had. That is the optimization-side echo of
Corollary~\ref{cor:irrelevance}.

\begin{remark}[The coordinate scale cancels]
\label{rem:scale}
The constant $c$ depends on how we chose to normalise $u$ and is not determined
by the task, so $\mu_0$ on its own is not parameter-free. But $\hat g$ is linear
in $c$, so $v_+$ scales as $c^2$, and the threshold \eqref{eq:scaling} sees the
two only through $\sqrt{v_+}/\mu_0$, where $c$ cancels. The numerical prediction
in Section~\ref{sec:experiments} is therefore immune to the choice of coordinate
scale. The budget criterion $\rho^{+}_{\mathrm{esc}} = a/(\mu_0S_T)$ is not,
because SGD is not invariant under reparameterisation of $u$; under AdamW, which
is scale-invariant in the gradient, both criteria are.
\end{remark}

\begin{lemma}[Per-example gradient variance]
\label{lem:variance}
Under \eqref{eq:gradform} and Assumption~\ref{ass:selection} at $u_0$, the
per-example gradient variance under $\cD_\rho$ satisfies
\begin{equation}
\label{eq:variance}
\sigma^2(\rho,\ell)/c^2
\;\le\;
\rho\,\bar A^\star(1-\bar A^\star)
+ (1-\rho)\,\tfrac1\ell\big(1-\tfrac1\ell\big)
+ \rho(1-\rho)\Delta^2/c^2 .
\end{equation}
Since the threshold sits at small $\rho$, the value that matters is the
$\rho\to0$ limit $v_+ = c^2\,\E[\tfrac1L(1-\tfrac1L)] \approx 0.1296\,c^2$.
\end{lemma}

\begin{proof}
At $u_0$ we have $p_{u_0}(J)=1/L$ deterministically, so within each component the
centred gradient is $c(\1\{J=I^\star\}-\pi)$ with $\pi=\Prob(J=I^\star)$: a
scaled centred Bernoulli of variance $c^2\pi(1-\pi)$. Proposition~\ref{prop:mu0}
gives $\pi = \bar A^\star$ on the good component and $1/\ell$ on the bad one.
Substitute into \eqref{eq:totalvar} and bound the between term.
\end{proof}

The $1/\ell$ that other treatments posit as a hypothesis is here a consequence of
$\Pi$ being uniform in Definition~\ref{def:traces}.

\subsection{What happens once the model moves}
\label{sec:bias}

Proposition~\ref{prop:mu0} holds at $u_0$, and the convenient cancellation there
does not survive elsewhere. From \eqref{eq:gooddrift},
\begin{equation}
\label{eq:bias}
\epsilon_\bad(u) \;=\; c\,\big|\E[p_u(J)] - \E[1/L]\big| ,
\end{equation}
which is $0$ at $u_0$ and generally positive away from it. So there is a feedback
loop lurking: correct traces push $u$ up, the model grows more confident in the
attended index, $\E[p_u(J)]$ rises above $\E[1/L]$, and now corrupted traces
start pulling back with force proportional to $(1-\rho)$.

\begin{proposition}[Interior threshold once bias is present]
\label{prop:interior}
If $\bar\epsilon > 0$, taking $B\to\infty$ in
$\rho\mu_0 - (1-\rho)\bar\epsilon \ge \kappa\sqrt{V/B}$ gives
\begin{equation}
\label{eq:interior}
\rho_c^\infty(\bar\epsilon) \;=\; \frac{\bar\epsilon}{\mu_0 + \bar\epsilon} > 0 .
\end{equation}
\end{proposition}

So the picture has two regimes: the threshold behaves as $\Theta(B^{-1/2})$ early,
near $u_0$, and drifts toward an interior value as $u_t$ moves away.

We stop short of bounding how far it drifts, and the reason is worth stating.
Doing so would need an assumption controlling $\E[p_u(J)]$ from above across all
of $\cU$, and a model that becomes overconfident---assigning the attended index
more probability than its accuracy warrants---breaks any such bound. Nothing in
the Bayes analysis prevents overconfidence, since $\E[p_u(J)]$ is a probability
and $\bar A^\star$ is an accuracy. So we make no claim of the form
$\rho_c^\infty\le\frac12$, tempting as it is.

\begin{remark}[Where the drift stalls, inside this coordinate]
\label{rem:stall}
By \eqref{eq:gooddrift} the expected trace-copy gradient vanishes when
$\E[p_u(J)] = \bar A^\star$. Read narrowly, this is a consistency check: the
optimization story does not predict progress past the ceiling the information
story establishes. It is a statement about one coordinate of one readout, not
about stationarity of the network or about its accuracy.
\end{remark}

\begin{remark}[A second axis we did not run]
\label{rem:adversarial}
You can make $\bar\epsilon$ positive by construction: interpolate between
$R^{\good}$ and $R^{\bad}$ so the corrupted trace points at a plausible but wrong
index instead of a uniform one. That gives an experimental axis orthogonal to
$\rho$, along which \eqref{eq:interior} predicts a monotonically rising interior
threshold. We did not implement it.
\end{remark}

% ==================================================================
\section{Objection 1: ``but you used Adam''}
\label{sec:adam}

Section~\ref{sec:scalar} is written for plain SGD and our runs use AdamW. Since
the headline empirical result is a mismatched exponent, the first thing to check
is whether the optimizer is the mismatch. This section builds a quasi-stationary
surrogate for AdamW on the same coordinate. The short answer is that adaptivity
changes the mechanism in three interesting ways and leaves the exponent alone.

We want to be clear about the status of what follows. The moment calculations
assume the gradient statistics are locally stationary over a window long compared
to the accumulator timescales, which at $\beta_2=0.999$ means $\gg10^3$ steps.
That is a demanding condition during warmup and across the transition itself,
where $u_t$ is moving fast. We use the surrogate to answer a comparative question
(does adaptivity change the exponent?), not to predict a trajectory.

\subsection{AdamW as three interacting variables}

\begin{definition}[AdamW on the trace-copy coordinate]
\label{def:adamw}
With $h_t := -\bar g_t$, AdamW iterates
\begin{align}
m_t &= \beta_1 m_{t-1} + (1-\beta_1)\,h_t, \label{eq:adam-m}\\
\nu_t &= \beta_2 \nu_{t-1} + (1-\beta_2)\,h_t^2, \label{eq:adam-v}\\
u_{t+1} &= u_t + \eta_t\,\frac{\hat m_t}{\sqrt{\hat\nu_t}+\varepsilon_{\mathrm A}}
           \;-\; \eta_t\lambda_\wdc\, u_t ,
\label{eq:adam-u}
\end{align}
with $\hat m_t = m_t/(1-\beta_1^t)$, $\hat\nu_t = \nu_t/(1-\beta_2^t)$,
$m_0 = \nu_0 = 0$, decoupled weight decay $\lambda_\wdc\ge0$, and numerical floor
$\varepsilon_{\mathrm A} > 0$. We write the floor as $\varepsilon_{\mathrm A}$ to
keep it clear of the bias $\epsilon_\bad$.
\end{definition}

One detail in \eqref{eq:adam-v} drives everything downstream: $\nu_t$ tracks the
second moment of the \emph{minibatch} gradient, not the per-example one. So it
carries its own factor of $1/B$, and the normalization in \eqref{eq:adam-u}
partially undoes the variance reduction that a larger batch just bought you.

\begin{assumption}[Quasi-stationary regime]
\label{ass:qs}
Fix $u\in\cU$ and a window of length
$T_w\gg\max\{(1-\beta_1)^{-1},(1-\beta_2)^{-1}\}$ over which
$\mu := \E[h_t\mid u]$ and $s^2 := \Var(h_t\mid u)$ are constant and $h_t$ is
independent across steps. Assume also a bounded normalized fourth moment,
$\Var(h_t^2)\le C_4(\mu^2+s^2)^2$, and that the accumulator covariance is
negligible at the order we retain,
$|\mathrm{Cov}(\hat m_t,\hat\nu_t)| = o\big(\mu\,\Var(\hat\nu_t)^{1/2}\big)$.
\end{assumption}

\begin{lemma}[Accumulator moments]
\label{lem:adammoments}
Under Assumption~\ref{ass:qs}, in the stationary regime
\begin{equation}
\label{eq:adammoments}
\E m_t = \mu,\quad
\Var(m_t) = \tfrac{1-\beta_1}{1+\beta_1}s^2,\quad
\E \nu_t = \mu^2 + s^2,\quad
\Var(\nu_t) = \tfrac{1-\beta_2}{1+\beta_2}\Var(h_t^2),
\end{equation}
and by \eqref{eq:batchvar}, $s^2\le V/B$.
\end{lemma}

\begin{proof}
Unroll: $m_t = (1-\beta_1)\sum_{k\ge0}\beta_1^kh_{t-k}$, so $\E m_t = \mu$ and
$\Var(m_t) = (1-\beta_1)^2s^2\sum_k\beta_1^{2k} = \frac{1-\beta_1}{1+\beta_1}s^2$.
The same computation with $h_t^2$ gives the $\nu$ statements, and bias correction
removes the $O(\beta_i^t)$ transient.
\end{proof}

\begin{remark}[Momentum is a bigger batch, and nothing more]
\label{rem:momentum}
Read $\Var(m_t)$ as noise from a batch of size $B(1+\beta_1)/(1-\beta_1)$. At
$\beta_1=0.9$ that is $19B$, better than a four-fold cut in gradient-noise
standard deviation. It is a large effect on the constant and no effect on the
exponent, since the factor does not involve $B$. The same one-liner disposes of
heavy-ball momentum as an explanation of our discrepancy.
\end{remark}

\subsection{The drift field: a saturating function of SNR}

\begin{proposition}[Saturating drift; quasi-stationary approximation]
\label{prop:adamdrift}
Under Assumption~\ref{ass:qs}, define
\begin{equation}
\label{eq:rdef}
r := \frac{\mu}{s} = \frac{\rho\mu_0\sqrt B}{\sqrt V},
\qquad
e_B := \frac{\varepsilon_{\mathrm A}}{s} = \frac{\varepsilon_{\mathrm A}\sqrt B}{\sqrt V}.
\end{equation}
Then the expected AdamW increment at $u$ is
\begin{equation}
\label{eq:adamphi}
\E\big[u_{t+1}-u_t \mid u_t = u\big]
= \eta_t\Big[\Phi(r,e_B) - \lambda_\wdc u\Big]
+ O\!\Big(\eta_t\sqrt{\tfrac{1-\beta_2}{1+\beta_2}}\Big),
\qquad
\Phi(r,e_B) := \frac{r}{\sqrt{1+r^2}+e_B},
\end{equation}
and the per-step fluctuation has standard deviation
$\eta_t c_1/(\sqrt{1+r^2}+e_B)$ with $c_1 := \sqrt{(1-\beta_1)/(1+\beta_1)}$.
\end{proposition}

\begin{proof}
Write $F(m,\nu) = m/(\sqrt\nu+\varepsilon_{\mathrm A})$ and expand about
$(\mu,\E\hat\nu_t)$. The zeroth-order term is
$\mu/(\sqrt{\mu^2+s^2}+\varepsilon_{\mathrm A})$. First-order terms vanish under
expectation, since $\E\hat m_t = \mu$ and $\E\hat\nu_t$ is the expansion point.
The leading survivors are the covariance term
$\partial^2_{m\nu}F\cdot\mathrm{Cov}(\hat m_t,\hat\nu_t)$, negligible by
Assumption~\ref{ass:qs}, and the curvature term
$\tfrac12\partial^2_{\nu}F\cdot\Var(\hat\nu_t)$. The bounded normalized fourth
moment gives $\Var(\hat\nu_t)/(\E\hat\nu_t)^2 = O((1-\beta_2)/(1+\beta_2))$, so
the relative fluctuation of $\sqrt{\hat\nu_t}$ is
$O(\sqrt{(1-\beta_2)/(1+\beta_2)})$, which is the order we retain. Divide
numerator and denominator by $s$ to get $\Phi(r,e_B)$; the decay term is exact.
The fluctuation claim follows from $\mathrm{sd}(\hat m_t) = c_1s$ and treating
the denominator as deterministic at the same order.
\end{proof}

We report the error at $O(\sqrt{(1-\beta_2)/(1+\beta_2)})$, matching the
denominator's fluctuation, rather than at the square of that rate. Claiming the
smaller error would require controlling the $\hat m$--$\hat\nu$ covariance, and
Assumption~\ref{ass:qs} assumes that away rather than establishing it. This is
the precise sense in which this section is a surrogate.

Three consequences deserve to be pulled out.

\paragraph{The drift is now bounded.}
Under SGD the drift is $\rho\mu_0$: linear in $\rho$, unbounded. Under AdamW it
is $\Phi(r,e_B)$, which behaves like $r/(1+e_B)$ for small $r$ and saturates at
$1/(1+e_B)$. The optimizer responds to the \emph{signal-to-noise ratio} of the
trace-copy gradient, not to its magnitude. Doubling the gradient does nothing if
you also doubled the noise.

\paragraph{Batch size enters through the opposite channel.}
Under SGD, per-step noise is $\eta_t\sqrt{V/B}$ and drift is $\eta_t\rho\mu_0$:
bigger batches quiet the noise. Under AdamW, per-step noise is at most
$\eta_tc_1$, independent of $B$, while the drift grows through $r\propto\sqrt B$:
bigger batches amplify the signal instead. The two theories move $B$ through
opposite doors. That the exponent survives the swap is the content of
Theorem~\ref{thm:adamthreshold}.

\paragraph{Weight decay removes the horizon from the criterion.}
\begin{proposition}[Weight decay converts a budget criterion into a threshold]
\label{prop:wd}
For $\lambda_\wdc > 0$ the deterministic flow
$\dot u = \Phi(r,e_B)-\lambda_\wdc u$ has a unique attracting fixed point
$u^\star = \Phi/\lambda_\wdc$. The mean trajectory therefore leaves $\cU$ in the
trace-reinforcing direction if and only if
\begin{equation}
\label{eq:wdcriterion}
\Phi(r,e_B) \;>\; \lambda_\wdc\,a ,
\end{equation}
and otherwise is confined for all time, not merely up to a horizon $T$.
\end{proposition}

This fixes something unsatisfying about Theorem~\ref{thm:escape}, which implies
that training longer should eventually wash out any threshold, whereas the curves
we measure are stable. Putting the two criteria together, the binding condition
is $\Phi(r,e_B)\ge\theta$ with
\begin{equation}
\label{eq:thetadef}
\theta := a\,\max\{\lambda_\wdc,\ S_T^{-1}\},
\end{equation}
the first branch being the weight-decay fixed point and the second the
displacement $S_T\Phi\ge a$ needed within the horizon. Both branches are free of
$B$, which is what lets the next theorem go through.

\begin{theorem}[Critical fraction under AdamW]
\label{thm:adamthreshold}
Under Assumption~\ref{ass:qs} in the zero-bias regime, with $\theta$ as in
\eqref{eq:thetadef} and $\theta(1+e_B) < 1$,
\begin{equation}
\label{eq:adamrhoc}
\rho_c^{\Adam}(B)
= \frac{\sqrt V}{\mu_0}\,\frac{\theta}{\sqrt{1-\theta^2}}\;B^{-1/2}
\;+\; \frac{\theta\,\varepsilon_{\mathrm A}}{\mu_0(1-\theta^2)}
\;+\; O\big(\varepsilon_{\mathrm A}^2\sqrt B\big),
\end{equation}
with local log--log slope
$-\tfrac12\big(1+e_B/\sqrt{1-\theta^2}\big)^{-1} \in (-\tfrac12,0]$.
\end{theorem}

\begin{proof}
Solve $F(r,e_B) := r-\theta\sqrt{1+r^2}-\theta e_B = 0$. At $e_B = 0$,
$r_0 = \theta/\sqrt{1-\theta^2}$. Since $\partial_rF = 1-\theta^2$ at $r_0$ and
$\partial_{e_B}F = -\theta$, the implicit function theorem gives
$dr/de_B = \theta/(1-\theta^2)$, so
$r_c = \theta/\sqrt{1-\theta^2}+\theta e_B/(1-\theta^2)+O(e_B^2)$. Substituting
$\rho_c = r_c\sqrt V/(\mu_0\sqrt B)$ and
$e_B = \varepsilon_{\mathrm A}\sqrt{B/V}$ makes the second term free of $B$.
Writing $\rho_c = c_1'B^{-1/2}+c_0'$ with
$c_0'/(c_1'B^{-1/2}) = e_B/\sqrt{1-\theta^2}$ gives the slope. It is negative and
above $-\tfrac12$ because $c_0'>0$: a $B$-independent additive floor can only
flatten a power law.
\end{proof}

\begin{remark}[The flattening is real and utterly invisible]
\label{rem:epsnumeric}
With $\varepsilon_{\mathrm A} = 10^{-8}$ and $\sqrt{v_+}/c = 0.360$, the
standardized floor at our largest batch is $e_{512}\approx6.3\times10^{-7}$. With
the calibrated $\theta = 0.985$ of Appendix~\ref{app:numbers}, the predicted
slope is $-0.5(1-3.7\times10^{-6})$. AdamW's numerical floor bends the law toward
\emph{shallower} decay by a few parts in a million, so it is the wrong sign and
the wrong size to explain anything we measured.
\end{remark}

\subsection{The same conclusion without AdamW}

Theorem~\ref{thm:adamthreshold} handles the optimizer we used. The next result
says the conclusion is not about AdamW in particular. We state the hypotheses
carefully, because the informal version (``scale-invariant optimizers can't change
the exponent'') is too strong to be true.

\begin{theorem}[Fixed-horizon exponent invariance]
\label{thm:noexp}
Consider $u_{t+1} = u_t + \eta_t\Psi(\bar g_{0:t}; u_{0:t})$ with $\Psi$
measurable and positively homogeneous of degree zero in its gradient arguments
jointly. Assume:
\begin{enumerate}[leftmargin=*,itemsep=1pt,label=(H\arabic*)]
\item the horizon $T$ is fixed as $B$ varies, and minibatches are drawn fresh, so
that conditionally on $\cF_t$ the batch gradient is a mean of $B$ i.i.d.\ draws;
\item the per-example gradient has a third absolute moment bounded uniformly over
$u\in\cU$, and $\Var(\hat g(u))$ is bounded away from $0$ on $\cU$;
\item the threshold is the level set of a criterion functional $\mathcal C$ on
the trajectory law that is free of $B$, continuous in the Gaussian limit, and
\emph{transverse} at the threshold: $\partial_r\mathcal C(r)\ne0$ at $r = r_c$.
\end{enumerate}
Then $\rho_c(B) = r_c\sqrt V\mu_0^{-1}B^{-1/2}(1+O(B^{-1/2}))$, so the asymptotic
log--log slope is $-1/2$.
\end{theorem}

\begin{proof}
Write $\bar g_t = -(\mu+s\zeta_t)$ with $\zeta_t$ standardized. By (H1) and (H2),
conditionally on $\cF_t$ the law of $\zeta_t$ is within Kolmogorov distance
$O(B^{-1/2})$ of $N(0,1)$ by Berry--Esseen, uniformly over $u\in\cU$; since $T$
is fixed, composing over steps perturbs the joint law of $(\zeta_t)_{t\le T}$ by
$O(TB^{-1/2}) = O(B^{-1/2})$. Degree-zero homogeneity applied with $c = 1/s$
gives $\Psi(\bar g_{0:t};u_{0:t}) =
\Psi(-(r+\zeta_0),\dots,-(r+\zeta_t);u_{0:t})$, so the trajectory is a fixed
measurable functional of $(r,(\zeta_t)_{t\le T})$. Hence $\mathcal C$ depends on
$(\rho,B)$ only through $r$, up to $O(B^{-1/2})$. Transversality (H3) then puts
the level set at $r = r_c + O(B^{-1/2})$, and substituting
$r = \rho\mu_0\sqrt B/\sqrt V$ gives the claim.
\end{proof}

\takeaway{Once the optimizer only sees the \emph{direction} of the gradient, the
pair $(\rho,B)$ can only reach it through the combination $\rho\sqrt B$. Any
threshold is then a level set of $\rho\sqrt B$, which is a $-1/2$ law.}

Hypothesis (H3) is doing real work: without transversality an $O(B^{-1/2})$
perturbation of the trajectory law can move a level set arbitrarily far.
Hypothesis (H1) is what our fixed-step protocol supplies and what a
fixed-sample-budget protocol would destroy, a point that returns in
Section~\ref{sec:gap}. Degree-zero homogeneity covers Adam, AdamW, RMSProp,
signSGD, Lion and normalized SGD. Plain SGD is degree one and is covered by
Section~\ref{sec:scalar}, where the two criteria of
Remark~\ref{rem:twothresholds} give slopes $-1/2$ and $0$.

\begin{remark}[Why $-1/2$ is hard to beat through the gradient law]
\label{rem:hardfloor}
For a mean of $B$ i.i.d.\ variables with finite variance,
$\Var(\bar g) = \Var(\hat g)/B$ exactly. No averaging scheme concentrates faster.
So any mechanism that acts through the minibatch-gradient law has $\rho$ and $B$
entering as $\rho\sqrt B$, and its thresholds are level sets of that combination.
The two candidate escapes have the wrong sign: heavy-tailed per-example gradients
with tail index $\alpha\in(1,2)$ concentrate at $B^{-(1-1/\alpha)}$, which is
\emph{shallower} than $B^{-1/2}$, and endogenous bias does the same by
Proposition~\ref{prop:biasflat}. This is a heuristic about the gradient law, not
a theorem about all mechanisms, and in particular it says nothing about channels
where $B$ enters from outside that law.
\end{remark}

\begin{proposition}[Bias feedback flattens]
\label{prop:biasflat}
Linearise \eqref{eq:bias} as $\epsilon_\bad(u) = \beta|u|$ with
$\beta := c\,\partial_u\E[p_u(J)]|_{u_0}\ge0$. Under AdamW in the small-$r$
regime, evaluating the net drift at the binding point $u = a$, criterion
\eqref{eq:wdcriterion} reads $(\rho\mu_0-(1-\rho)\beta a)\sqrt{B/V}\ge\theta$,
whence
\begin{equation}
\label{eq:biasflat}
\rho_c(B) = \frac{\beta a + \theta a\sqrt{V/B}}{\mu_0 + \beta a},
\qquad
\frac{d\log\rho_c}{d\log B}\in\Big(-\tfrac12,\ 0\Big].
\end{equation}
\end{proposition}

That settles the sign of the bias mechanism: it raises the threshold at every $B$
and pushes the exponent toward $0$, consistent with
Proposition~\ref{prop:interior}. The horizon behaves the same way, since by
\eqref{eq:thetadef} it enters AdamW only inside the $B$-free constant $\theta$.

% ==================================================================
\section{Objection 2: ``but your model has 810{,}000 parameters''}
\label{sec:matrix}

The second obvious complaint is that a one-dimensional analysis has no business
speaking for a real network. This section lifts escape and persistence to the
parameter manifold of a \emph{linear} attention layer, and finds that dimension
shows up in the persistence half of the argument and not in the threshold.

The scope, stated plainly: this is a surrogate. It shows that dimensionality by
itself does not change the leading exponent in a softmax-free attention layer
that satisfies explicit alignment assumptions. It does not remove the
scalar-reduction objection for the four-layer softmax transformer we actually
train.

\begin{definition}[Linear attention block]
\label{def:linattn}
Let $X\in\R^{n\times d}$ be the residual stream and $x_q$ the answer-position
row. A linear attention head computes
$\mathrm{Attn}(X) = W_{OV}X^\top XW_{QK}^\top x_q
= \sum_j\langle W_{QK}^\top x_q, x_j\rangle W_{OV}x_j$, and the answer logits are
$\ell = E^\top\mathrm{Attn}(X)$ for a fixed unembedding $E$. Collect the
trainable parameters into $U := (W_{QK},W_{OV})\in\R^D$ with $D = 2d^2$, and
write $\cU_a := \{U : \|U-U_0\|_F < a\}$.
\end{definition}

\begin{definition}[Trace-copy direction]
\label{def:tracedir}
With $G^\good(U) := \E[-\nabla_U\cL\mid Z=1]$, set
$\Om^\star := G^\good(U_0)/\|G^\good(U_0)\|_F$ and
$u_t := \langle U_t - U_0,\Om^\star\rangle$.
\end{definition}

\begin{assumption}[Trace-copy alignment]
\label{ass:alignment}
There is $\delta\ge0$ with
$\langle\nabla_U\ell_k,\Om^\star\rangle = c\,\1\{k = J\}+O(\delta)$ on $\cU_a$,
$c>0$: moving along $\Om^\star$ boosts the copied index selectively and is
$\delta$-incoherent with the other logit directions.
\end{assumption}

Normalising the good-component gradient fixes a direction, but it does not by
itself make that direction selective for the copied index, so
Assumption~\ref{ass:alignment} is an assumption and not a consequence of
Definition~\ref{def:tracedir}.

\begin{proposition}[The scalar gradient form is a projection]
\label{prop:projection}
The per-example matrix gradient is
$-\nabla_U\cL = \sum_k(\1\{k=I^\star\}-p_U(k))\nabla_U\ell_k$, so under
Assumption~\ref{ass:alignment}
\begin{equation}
\label{eq:matproj}
\langle-\nabla_U\cL,\ \Om^\star\rangle
= c\big(\1\{J=I^\star\}-p_U(J)\big) + O(\delta),
\end{equation}
which is \eqref{eq:gradform} up to the remainder. So
Proposition~\ref{prop:mu0}, Lemma~\ref{lem:variance} and \eqref{eq:bias} all
transfer to the projected coordinate, with the same constant $c$.
\end{proposition}

\begin{proof}
The displayed gradient is the softmax--cross-entropy residual in vector form.
Take the inner product with $\Om^\star$ and apply
Assumption~\ref{ass:alignment} termwise: the $k=J$ term contributes
$c(\1\{J=I^\star\}-p_U(J))$, and the rest contribute $O(\delta)$ each against
coefficients whose absolute values sum to at most $2$.
\end{proof}

\begin{assumption}[Matrix drift, covariance and tails]
\label{ass:matrix}
On $\cU_a$: (1) $\E[-\hat G^\good(U)] = \mu(U)\Om^\star+\Xi^\good(U)$ with
$\mu_0\le\mu\le\mu_1$ and $\|\Xi^\good\|_F\le\bar\epsilon_\perp$;
(2) $\E[-\hat G^\bad(U)] = \Xi^\bad(U)$ with $\|\Xi^\bad\|_F\le\epsilon_\bad(U)$;
(3) $\mathrm{Cov}(\vecop\hat G)\preceq\Sigma$, with
$v_\parallel := \langle\Om^\star,\Sigma\Om^\star\rangle\le v_+$ and
$D_{\mathrm{eff}} := \Tr\Sigma/v_+$;
(4) $\langle\hat G,\Om\rangle$ is sub-exponential with parameter $K_e$ for every
unit $\Om$.
\end{assumption}

Note that $D_{\mathrm{eff}}$, not $D$, is what appears: by \eqref{eq:matproj}
each per-example gradient is a sum of a few rank-one features, so $\Sigma$ has a
rapidly decaying spectrum.

\begin{theorem}[Matrix escape and persistence]
\label{thm:matrixescape}
Let Assumption~\ref{ass:matrix} hold with
$\epsilon_\bad\equiv0$, $\bar\epsilon_\perp = 0$, and write
$V_\parallel := v_\parallel+\rho(1-\rho)\mu_1^2$ and
$V_\perp := \Tr\Sigma+\rho(1-\rho)\mu_1^2$. Then:
\begin{enumerate}[leftmargin=*,itemsep=2pt]
\item[(i)] if $\rho > a/(\mu_0S_T)$, the trajectory exits $\cU_a$ in the
trace-reinforcing direction by time $T$ with probability at least
$1-\exp(-\lambda^2/[2(C_\eta V_\parallel/B+K_e\lambda/3)])$ with
$\lambda := \rho\mu_0S_T-a$, and \emph{this bound does not involve $D$};
\item[(ii)] if $\rho < a/(\mu_1S_T)$, the trajectory stays in $\cU_a$ for all
$t\le T$ with probability at least
$1-(D+1)\exp(-\lambda'^2/[2(C_\eta V_\perp/B+K_e\lambda'/3)])$ with
$\lambda' := a-\rho\mu_1S_T$.
\end{enumerate}
\end{theorem}

Part (i) is proved by projecting onto $\Om^\star$, after which the process is the
scalar recursion with $V\mapsto V_\parallel$; no union bound over directions is
needed, which is why $D$ never appears. Part (ii) has to control all $D$
coordinates at once, and is proved with the matrix Freedman inequality of
\citet{tropp2011freedman} applied to the self-adjoint dilation of the
$\R^D$-valued martingale. Appendix~\ref{app:matrix} gives both arguments,
including the dilation and the variance-process bound. The dimensional factor
there is the standard one for that inequality and we have not tried to optimise
it.

\begin{corollary}[Dimension-free exponent]
\label{cor:dimfree}
$\rho_c^{\mathrm{mat}}(B,\kappa) = \kappa\sqrt{v_\parallel}\mu_0^{-1}B^{-1/2}
+O(B^{-1}) = \Theta(B^{-1/2})$, uniformly in $D$. Dimension enters only the
persistence half, widening the undecided window by $\sqrt{D_{\mathrm{eff}}}$ and
inflating its failure probability by $D+1$.
\end{corollary}

\begin{remark}[Why spreading the signal does not dilute it]
\label{rem:spread}
The natural intuition is that distributing the signal over $D$ coordinates should
weaken it by $1/\sqrt D$. It does not, because the noise gets projected along
with the signal. The drift along $\Om^\star$ is $\|G^\good\|_F$ no matter how
many coordinates carry it, while the noise along $\Om^\star$ is a variance in
\emph{one} direction. The ratio that governs escape is intrinsically
one-dimensional, which is the real reason the scalar reduction is not as naive as
it looks.
\end{remark}

Coordinate-wise adaptivity does not interact with the geometry in a way that
changes this either. Applying Proposition~\ref{prop:adamdrift} coordinate-wise
and projecting gives drift $\eta_t\sum_i\Om^\star_i\Phi(r_i,e_{B,i})$, which in
the small-$r$ regime is
$\eta_t\sqrt B\rho\mu_0\sum_i\Om^{\star2}_i/\sqrt{\Sigma_{ii}}\,(1+o(1))$: once
again a $B$-free level set of $\rho\sqrt B$.

\begin{remark}[One initialization where the law would differ]
\label{rem:saddle}
Suppose the readout were degenerate at initialization, $\mu(U_0) = 0$, with drift
linear in displacement along the unstable direction of the bilinear map. Then
escape would be seeded by noise rather than driven by drift, and
$\rho_c\asymp(\log B+C)/(2\mu_0'S_T)$, which \emph{increases} with $B$: larger
batches supply less of the noise a saddle escape needs. Our task is not in that
regime, since Proposition~\ref{prop:mu0} gives $\mu_0>0$ at the symmetric
initialization. We mention it because it shows the exponent is not a universal
fact about batch size; it is a fact about this mechanism.
\end{remark}

% ==================================================================
\section{Experiments}
\label{sec:experiments}

\subsection{Setup}

We train a $4$-layer, $4$-head, width-$128$ causal transformer
($\approx8.1\times10^5$ parameters) with a character-level vocabulary and context
length $64$. Optimizer AdamW, learning rate $3\times10^{-4}$ with $600$ warmup
steps and cosine decay to $10^{-5}$, weight decay $0.01$, dropout $0.05$, label
smoothing $0.1$, gradient clipping at $1.0$, $6000$ steps per run. Every
$(\rho,B)$ configuration is run at three seeds $(2200,1337,2026)$ controlling
initialization and batch order, giving $165$ runs.

The training pool holds $1{,}536{,}000$ examples, generated once and shared
across all runs. For a given $\rho$, the first $\rho N$ examples carry correct
traces and the rest corrupted ones, and the correct and corrupted variants of the
same instance are never both used.

\paragraph{What ``accuracy'' means here, exactly.}
This matters more than it usually does, so we spell it out. Evaluation uses $500$
held-out prompts. The model gets the \emph{bare prompt} $W\texttt{;}C$---the
trace is \emph{not} supplied---and generates greedily for up to $35$ new tokens
or until a newline. We then parse the answer by taking the segment after the
final colon and reading the last digit string in it, and score exact string
match.

That is a compound, end-to-end metric, not teacher-forced accuracy at the answer
position. It bundles together emitting a well-formed trace, reaching the
delimiter, predicting the answer, and staying inside the generation budget.
Section~\ref{sec:anomaly} shows the bundling is not hypothetical: at $\ell=10$
the answer lands on token $34$ of a $35$-token budget, so one spurious token
truncates it. We report the metric we computed and flag the teacher-forced
version as the primary fix in Section~\ref{sec:limitations}.

\subsection{The transition, read against computed reference levels}

Table~\ref{tab:batchsweep} gives accuracy against $\rho$ at five batch sizes over
three seeds. Read any column top to bottom. The low-$\rho$ baseline sits at
$15.3$--$17.0\%$ against a computed floor of $15.7\%$, and the plateau at
$87$--$90\%$ against a computed ceiling of $89.3\%$. Nothing was fitted to
produce those two numbers, and away from the transition the seed-to-seed standard
deviations are under $2.5$ points, so this is not seed selection.

A pleasant side effect: because the metric is end-to-end, reaching the Bayes
ceiling means trace generation and parsing cost essentially nothing in the
plateau regime. Which makes their apparent cost at $\rho=1$ all the more
conspicuous.

The transition itself is sharp. At $B=128$, accuracy goes from $27.2\%$ at
$\rho=0.20$ to $81.6\%$ at $\rho=0.40$, which in normalized terms is $16\%$ to
$90\%$ of the available range, over a single doubling of $\rho$.

\paragraph{Seed variance peaks inside the transition, which is the point.}
The standard deviations are strongly $\rho$-dependent. At $B=256$ they run
$1.0, 0.4, 5.8, 18.3, 24.9, 1.9, 1.7,\dots$ as $\rho$ increases, peaking at
$\rho=0.20$ and collapsing on both sides. That $24.9$ comes from three runs at
$43.9\%$, $88.1\%$ and $46.2\%$: one seed escaped, two did not.

This is the most direct evidence in the paper for the escape-versus-confinement
picture, as against a story where the circuit simply improves smoothly. Near
$r\approx1$ the escape event of Theorem~\ref{thm:escape} is genuinely stochastic,
so what you should see at intermediate $\rho$ is a coin flip between two outcomes
rather than a smooth intermediate accuracy. That is what we see.

\begin{table}[t]
\caption{Accuracy (\%) versus correct-trace fraction $\rho$ at five batch sizes,
with steps fixed at $6000$. Mean $\pm$ s.d.\ over three seeds. Chance is $15.7\%$
and the Bayes ceiling $89.3\%$; neither is fitted. Bold marks cells inside each
column's transition. $\hat\rho_c$ is the logistic fit to the seed-averaged curve
with the $\rho = 1$ row excluded.}
\label{tab:batchsweep}
\centering
\small
\begin{tabular}{lccccc}
\toprule
$\rho$ & $B=32$ & $B=64$ & $B=128$ & $B=256$ & $B=512$ \\
\midrule
0.00 & $16.0_{\pm1.8}$ & $16.3_{\pm2.1}$ & $17.0_{\pm2.3}$ & $16.1_{\pm1.0}$ & $16.4_{\pm0.9}$ \\
0.05 & $16.0_{\pm0.9}$ & $15.3_{\pm1.4}$ & $16.1_{\pm0.2}$ & $16.6_{\pm0.4}$ & $16.7_{\pm1.2}$ \\
0.10 & $16.1_{\pm0.4}$ & $16.4_{\pm0.6}$ & $17.8_{\pm2.0}$ & $20.0_{\pm5.8}$ & $\mathbf{54.3}_{\pm29.6}$ \\
0.15 & $16.3_{\pm1.4}$ & $16.3_{\pm1.4}$ & $21.4_{\pm3.0}$ & $\mathbf{48.4}_{\pm18.3}$ & $\mathbf{85.7}_{\pm7.2}$ \\
0.20 & $16.7_{\pm1.6}$ & $18.0_{\pm1.5}$ & $\mathbf{27.2}_{\pm9.9}$ & $\mathbf{59.4}_{\pm24.9}$ & $88.8_{\pm1.4}$ \\
0.30 & $16.8_{\pm0.2}$ & $20.8_{\pm4.4}$ & $\mathbf{47.9}_{\pm11.6}$ & $88.0_{\pm1.9}$ & $88.0_{\pm2.0}$ \\
0.40 & $19.7_{\pm2.4}$ & $\mathbf{41.2}_{\pm15.0}$ & $\mathbf{81.6}_{\pm11.5}$ & $88.4_{\pm1.7}$ & $90.2_{\pm0.8}$ \\
0.50 & $\mathbf{27.9}_{\pm10.0}$ & $\mathbf{61.7}_{\pm23.5}$ & $87.9_{\pm2.3}$ & $89.1_{\pm1.0}$ & $89.7_{\pm0.8}$ \\
0.60 & $\mathbf{35.7}_{\pm13.1}$ & $\mathbf{73.6}_{\pm8.3}$ & $89.6_{\pm0.1}$ & $89.1_{\pm0.6}$ & $89.3_{\pm0.7}$ \\
0.80 & $\mathbf{56.7}_{\pm11.2}$ & $87.1_{\pm3.3}$ & $87.9_{\pm2.8}$ & $88.1_{\pm2.0}$ & $89.2_{\pm1.0}$ \\
1.00 & $39.5_{\pm14.9}$ & $59.8_{\pm12.7}$ & $61.7_{\pm16.7}$ & $71.7_{\pm3.3}$ & $79.2_{\pm6.5}$ \\
\midrule
$\hat\rho_c$ & 0.759 & 0.465 & 0.305 & 0.175 & 0.099 \\
\bottomrule
\end{tabular}
\end{table}

\subsection{The threshold does vanish with batch size, at a rate we did not predict}
\label{sec:exponent}

To extract a threshold we fit the logistic profile
\[
A(\rho) = A_{\mathrm{chance}} + \frac{\bar A^\star-A_{\mathrm{chance}}}
{1+\exp(-(\rho-\rho_c)/\tau)},
\]
with both levels pinned at their computed values, leaving only $\rho_c$ and the
width $\tau$ free, and excluding the $\rho=1$ row.

\paragraph{Good news.}
$\hat\rho_c$ falls monotonically from $0.759$ at $B=32$ to $0.099$ at $B=512$: a
factor of $7.7$ across a $16$-fold increase, and monotone within every seed
separately. There is no batch-size-independent critical trace fraction, and any
theory predicting a fixed interior threshold is refuted by these data.

\paragraph{Bad news.}
Regressing $\log\hat\rho_c$ on $\log B$ over all fifteen $(B,\mathrm{seed})$
cells gives
\begin{equation}
\label{eq:measured}
\hat\rho_c \propto B^{-0.72\pm0.05},\qquad R^2 = 0.95,
\end{equation}
against a predicted $-1/2$. The seed-averaged curve gives $-0.73$ with
$R^2 = 0.997$, and a midpoint-crossing estimator gives $-0.74$, so this is not an
artefact of the fitting profile. Table~\ref{tab:perseed} has the per-seed fits
and Figure~\ref{fig:phase} plots the curves against the two computed levels.

\paragraph{How firmly is $-1/2$ actually excluded?}
Two defensible analyses of the same three seeds disagree, and we would rather
show both than pick one quietly.

Treat each of the fifteen fitted thresholds as a data point and you get
$-0.72\pm0.05$, a $95\%$ interval of $[-0.80,-0.64]$, excluding $-1/2$.

Treat the \emph{seed} as the unit of replication---the conservative reading,
since all runs share one data pool---and the three per-seed exponents are
$-0.794$, $-0.782$ and $-0.570$, with mean $-0.715$ and standard error $0.073$.
On two degrees of freedom the $95\%$ interval is $[-1.03,-0.40]$, which contains
$-1/2$, at $p=0.10$.

We report the first as the point estimate and the second as the honest
uncertainty. What can be said without qualification is that all three seeds are
steeper than $-1/2$, and that the third seed's $-0.570$ hangs on a single cell:
$\hat\rho_c = 0.131$ at $B=512$ against $0.089$ for the other two.

\begin{table}[t]
\caption{Per-seed thresholds and exponents. Every seed gives a monotonically
falling $\hat\rho_c$ and an exponent steeper than $-1/2$, but the spread is large
enough that a seed-level test does not exclude $-1/2$.}
\label{tab:perseed}
\centering
\small
\begin{tabular}{lccccccc}
\toprule
Seed & $B{=}32$ & $B{=}64$ & $B{=}128$ & $B{=}256$ & $B{=}512$ & exponent & $R^2$ \\
\midrule
2200 & 0.849 & 0.557 & 0.349 & 0.208 & 0.088 & $-0.794$ & 0.977 \\
1337 & 0.754 & 0.415 & 0.266 & 0.132 & 0.089 & $-0.782$ & 0.994 \\
2026 & 0.649 & 0.420 & 0.301 & 0.198 & 0.131 & $-0.570$ & 0.999 \\
\midrule
mean curve & 0.759 & 0.465 & 0.305 & 0.175 & 0.099 & $-0.730$ & 0.997 \\
\bottomrule
\end{tabular}
\end{table}

\begin{figure}[t]
\centering
\includegraphics[width=0.72\textwidth]{figures/phase_transition_low_variance.png}
\caption{Accuracy versus correct-trace fraction $\rho$. The horizontal lines are
the computed chance floor $0.157$ and Bayes ceiling $0.893$ of
Corollary~\ref{cor:levels}; neither is fitted to the data.}
\label{fig:phase}
\end{figure}

\subsection{A surprise at \texorpdfstring{$\rho=1$}{rho=1}}
\label{sec:anomaly}

At every batch size, training on $100\%$ correct traces is worse than training on
$80\%$: $56.7\to39.5$, $87.1\to59.8$, $87.9\to61.7$, $88.1\to71.7$,
$89.2\to79.2$. Nothing in this paper predicts that. The AdamW drift $\Phi(r)$
increases with $\rho$ at every $B$, and so does the supervision-density account
of Section~\ref{sec:competing}.

\paragraph{It is not run-to-run noise. It could still be a data artefact.}
The effect appears in all fifteen cells individually, with per-cell deficits from
$3.4$ to $47.0$ points, so it survives three independent initializations and
batch orderings. That does not make it intrinsic. Every run draws from the same
pool, and at $\rho=1$ that pool is used in its entirety with correct traces, so
the \emph{data are identical} across all fifteen runs and only initialization and
batch order vary. Regenerating the pool under a different generator seed is the
control we are missing.

\paragraph{Our leading suspect is our own evaluation script.}
Recall the $35$-token generation cap. For a word of length $\ell$ the model must
emit a trace of $3\ell-1$ characters, then \texttt{ : }, then the index: $3\ell+3$
tokens before the answer, plus a newline. At $\ell=10$ that is exactly $34$ of
the $35$ available. At $\ell=9$ it is $31$.

So a single spurious token wipes out the answer entirely on the longest words,
which are also the words with the lowest ceiling. Any effect that makes a model
emit longer or less well-calibrated traces will register as an accuracy loss
under this metric even when the answer-position representation is perfectly
fine. And $\rho=1$ is exactly the regime where the model has never seen an
unpredictable trace, is maximally committed to generating a full one, and the
budget binds hardest.

This is testable and cheap: under teacher-forced answer scoring, or with a larger
token budget, the $\rho=1$ deficit should shrink or disappear.

\paragraph{If the confound is ruled out, two substantive candidates remain.}
First, \emph{regularisation by corrupted traces}: a minority of corrupted traces
may stop the model over-relying on a brittle copy circuit, forcing it to maintain
the prompt-only pathway that actually achieves Bayes accuracy. Second,
\emph{optimization pathology at the boundary}: at $\rho=1$ the trace is perfectly
predictable, trace-token loss collapses quickly, and the effective gradient
reaching the answer token may be reshaped. The answer-only ablation of
Section~\ref{sec:competing} discriminates the second.

One caution about reading too much into the shape. The deficits, $17.1, 27.3,
26.2, 16.4, 10.0$, are non-monotone in $B$ and peak at $B=64$--$128$. An earlier
single-seed version of this table looked like a clean monotone decay; that trend
did not survive replication.

% ==================================================================
\section{Chasing the missing $0.22$}
\label{sec:gap}

Time to collect what the theory has to say about the $-0.72$ measurement. The
summary: the two natural objections do not explain it, two further mechanisms
have the wrong sign, and what is left is sample exposure plus the possibility
that three seeds cannot tell $-0.72$ from $-0.5$.

\paragraph{Optimizer: no.}
Section~\ref{sec:adam} carries the argument through for AdamW under
Assumption~\ref{ass:qs} and lands on $-1/2$. The numerical-floor correction points
the wrong way and is of relative size $4\times10^{-6}$
(Remark~\ref{rem:epsnumeric}). Theorem~\ref{thm:noexp} extends the conclusion to
scale-invariant rules at fixed horizon under transversality, and momentum is a
$B$-free constant factor.

\paragraph{Dimension: no.}
Corollary~\ref{cor:dimfree} gives $\Theta(B^{-1/2})$ uniformly in $D$ in the
linear-attention surrogate, subject to that surrogate's scope.

\paragraph{Endogenous bias and horizon: wrong sign.}
Proposition~\ref{prop:biasflat} shows bias feedback adds a $B$-independent term
to $\rho_c$, which moves the slope toward $0$, not past $-1/2$. Under AdamW the
horizon enters only through the $B$-free $\theta$ of \eqref{eq:thetadef}.

\paragraph{Sample exposure: the one we cannot dismiss.}
At fixed steps, raising $B$ does two things at once. It raises the
signal-to-noise ratio, which is the channel our theory models, and it raises the
total number of examples seen, $N = BT$, which is a channel entirely outside the
minibatch-gradient law. Neither Theorem~\ref{thm:noexp} nor
Remark~\ref{rem:hardfloor} covers it, because both are statements about how $B$
enters through the gradient.

The following identity converts our measured exponent into a claim about that
channel, which is at least a quantity someone could go and measure.

\begin{proposition}[Sample-budget susceptibility]
\label{prop:susceptibility}
Suppose the measured threshold is the level set of a criterion of product form
$\Phi(r)\Lambda(N)\ge\theta$, where $\Lambda$ is increasing and represents
progress on a $\rho$-independent component of the task that gates the measured
accuracy: learning the output format, the prompt-search pathway, or reliable
trace generation. Then at fixed $T$,
\begin{equation}
\label{eq:susceptibility}
\frac{d\log\rho_c}{d\log B} = -\frac12 - \chi,
\qquad
\chi := \frac{d\log\Lambda/d\log N}{d\log\Phi/d\log r}\ \ge 0 .
\end{equation}
In the small-$r$ regime $d\log\Phi/d\log r\to1$, so matching $-0.72$ requires
$d\log\Lambda/d\log N\approx0.22$.
\end{proposition}

\begin{proof}
Take logarithms at equality and differentiate with respect to $\log B$ at fixed
$T$, using $\log N = \log B+\log T$ and
$\log r = \log\rho+\tfrac12\log B+\mathrm{const}$.
\end{proof}

A sample-budget exponent of $0.22$ sits comfortably inside the range reported for
data-scaling laws. We offer that as a plausibility check, not as evidence.
Propagating the fifteen-point standard error gives $\chi = 0.22\pm0.05$;
propagating the seed-level error gives $\chi = 0.22\pm0.07$, an interval that
contains $0$. The channel is available and quantified. It is not established.

\paragraph{Why the obvious control does not do what it looks like it does.}
The tempting fix is to hold $N = BT$ fixed and vary $B$, on the grounds that this
removes the sample-exposure channel. It does remove that channel, and it does not
isolate gradient noise. Fixing $N$ sets $T = N/B$, so the horizon now changes
with $B$. That violates (H1) of Theorem~\ref{thm:noexp} and drags back in exactly
the budget confound of Remark~\ref{rem:twothresholds}, with $S_T$ changing too.
Fixed-step, fixed-sample and fixed-epoch protocols identify three different
compute allocations, and none of them isolates minibatch-noise scaling on its
own. There is no free lunch here; the way out is to move one variable that is not
$B$.

\begin{prediction}[Data-diversity control]
\label{pred:diversity}
Fix $B = 256$ and $T = 6000$, and vary only the number of \emph{distinct}
examples in the training pool over $\{N/16,\dots,N\}$, sampling with repetition
so that batch size, step count and cumulative step size are all unchanged. Then
$r$ is unchanged by construction and \eqref{eq:susceptibility} collapses to
$d\log\rho_c/d\log N_{\mathrm{pool}} = -\chi$, so the sweep measures $\chi$
directly. Under the sample-exposure account it should return $\chi\approx0.22$.
Under the null---the theory is right and the measured exponent is seed noise---it
should return $\chi\approx0$.
\end{prediction}

We prefer this to the fixed-sample-budget sweep because it moves one quantity and
one only. The caveat is that repetition changes the effective data distribution,
so the control is clean only if the pool stays large enough to avoid memorisation
at the smallest setting.

% ==================================================================
\section{The rival explanation, and the experiment that separates them}
\label{sec:competing}

There is a second story for a $\rho$-dependent transition that has nothing to do
with trace-copy drift, and the standard training convention hands it to us for
free.

Under \eqref{eq:lfull} the loss is averaged over trace and answer tokens alike.
In our task $|R|\approx3\ell$ and $|Y|\approx1$, so about $95\%$ of the
supervised tokens are trace tokens. When traces are correct, modelling them
shares structure with answering, so the effort partly transfers. When they are
corrupted they are unpredictable by construction, so the optimizer burns most of
its budget fitting noise and the answer-token signal is diluted. Fitting
incompressible noise on most supervised positions may also degrade the
representations that the answer position depends on.

Under this account the transition happens where the useful fraction of the
gradient budget crosses a threshold, with no trace-copy circuit anywhere in the
story. Note that this account also predicts a batch-size dependence, so the mere
existence of a batch-size effect does not tell the two apart.

They do differ in one testable respect.

\begin{definition}[Supervision conditions]
\label{def:conditions}
\textbf{A, full supervision}: $\cL_{\mathrm{full}}$ of \eqref{eq:lfull}.
\textbf{B, answer-only}: the model ingests prompt \emph{and} trace, but
cross-entropy is masked to the answer tokens,
$\cL_{\mathrm{ans}} = |Y|^{-1}\sum_{j\in Y}\cL_{\mathrm{CE}}(y_j,y_j^\star)$.
\textbf{C, direct answer}: no trace in the sequence at all, $W\to Y$.
\end{definition}

Conditions A and B present identical token sequences and differ only in the loss
mask; in our implementation the supervised-token fraction is $41.2\%$ under A and
$4.8\%$ under B. Condition C has no $\rho$ axis at all: once the trace is
removed, a correct-trace and a corrupted-trace example are byte-identical, so C
is one baseline number per configuration and belongs on the plot as a horizontal
line.

Here is the discrimination. Under the trace-copy account, the mechanism is drift
from the answer-token gradient \eqref{eq:gradform}, which is present in both A
and B, so the transition should persist under B at a comparable $\rho_c$. Under
the supervision-density account, the mechanism is dilution by trace-token loss,
which B removes outright, so the transition should shift a long way or vanish. We
adopt a shift tolerance of $0.10$ in $\rho_c$.

One guard is necessary for either account to be testable at all: we refuse to
report a threshold when the curve rises by less than $15\%$ of the
chance-to-ceiling span, or when the fitted $\rho_c$ lands outside $[-0.25,1.25]$.
This is not pedantry. An undertrained run produces a flat curve, and a logistic
fit will happily return a confident and meaningless $\rho_c$ from it.

Condition B has a bonus given Section~\ref{sec:anomaly}: because the trace is
supplied rather than generated, evaluating under B with the trace in context also
removes the generation-length confound.

% ==================================================================
\section{Limitations}
\label{sec:limitations}

Listed worst first.

\paragraph{The evaluation metric is end-to-end, not answer-position.}
Accuracy comes from free generation off the bare prompt followed by regex
parsing, under a $35$-token cap that binds at $\ell=10$. That confounds trace
generation, delimiter emission, answer prediction and parsing. It affects the
$\rho=1$ anomaly directly and, in principle, every number in
Table~\ref{tab:batchsweep}. The fix is to report teacher-forced accuracy at the
answer position with a clean trace supplied in context, which is the quantity the
theory is actually about. We have not run it. That the plateau lands on the
computed Bayes ceiling does bound how much the confound can be distorting the
high-$\rho$ regime, but says nothing about low $\rho$ or about $\rho=1$.

\paragraph{Three seeds is enough to see the effect, not to bound it.}
Three seeds settled two questions: the $\rho=1$ anomaly is not run-to-run noise,
and seed variance does peak inside the transition as predicted. A third question
is still open, since the fifteen-point interval on the exponent excludes $-1/2$
and the seed-level interval does not. Only more seeds will separate them. All
seeds also share one data pool, so pool-specific artefacts remain unexcluded
everywhere in this paper.

\paragraph{The AdamW analysis is a quasi-stationary surrogate.}
Assumption~\ref{ass:qs} wants locally constant gradient statistics over a window
long compared to $(1-\beta_2)^{-1}\approx10^3$ steps, plus a bounded fourth
moment and negligible accumulator covariance. Neither the stationarity nor the
covariance condition holds exactly during warmup or across the transition. A full
treatment would track the joint law of $(u_t,m_t,\nu_t)$. Note that
Theorem~\ref{thm:noexp} does not lean on quasi-stationarity---only on degree-zero
homogeneity, a fixed horizon, uniform third moments and transversality---so the
invariance conclusion survives even if the surrogate's constants do not.

\paragraph{The matrix analysis is a linear-attention surrogate.}
Definition~\ref{def:linattn} drops the softmax, and
Assumptions~\ref{ass:alignment} and~\ref{ass:matrix} are assumptions rather than
consequences. What is established is that dimensionality alone does not change
the leading exponent in that surrogate. What is not established is that the
scalar reduction is safe for the four-layer softmax transformer we train.

\paragraph{The segment-selection assumptions are substantive.}
Assumption~\ref{ass:selection} is what makes $\Prob(J=I^\star) = 1/K$ on correct
traces and $1/L$ on corrupted ones. Its second clause restricts what a trained
model may do: it could select segments using the corrupted characters or the
positional structure, and then $\mu_0$ and the bias both change. The assumptions
are most plausible near initialization, which is where we use them, and least
plausible late in training, which is where we do not.

\paragraph{Per-example independence of $Z$ is a scope condition, not a detail.}
Section~\ref{sec:variance} shows that a shared correctness indicator within a
minibatch would leave the between-component variance outside the $1/B$, and the
threshold would then not vanish at all. Our prediction applies to i.i.d.\ mixed
pools. It does not apply to regimes where correct and corrupted data arrive in
homogeneous blocks, which is arguably the more realistic setting for distilled
pipelines.

\paragraph{This task has no information gap, by construction.}
Per Remark~\ref{rem:gap}, everything here concerns a setting where the trace is
provably useless at the Bayes level. Whether the same picture holds when the
trace genuinely carries information is open, and it is the variant we would build
next.

% ==================================================================
\section{Conclusion}
\label{sec:conclusion}

We set out to explain why accuracy depends so sharply on the fraction of correct
think traces. The first result was negative and clarifying: in this task the
trace is conditionally independent of the answer given the prompt, so the
transition cannot be about information. Both the ceiling and the floor are
closed-form, and both are hit by experiment with nothing fitted.

That relocated the phenomenon to the optimizer. There, the critical trace
fraction comes from a competition between a drift proportional to $\rho$ and
noise proportional to $B^{-1/2}$, with the drift constant supplied by the task's
own combinatorics and the leftover coordinate scale cancelling out. The
prediction is that the critical fraction vanishes as $B^{-1/2}$.

Experiment confirms the vanishing---a factor of $7.7$ over a $16$-fold range of
$B$---and measures $-0.72\pm0.05$, with every seed steeper than $-1/2$ but a
conservative seed-level interval that still contains it. We then chased the two
natural objections. Under a quasi-stationary surrogate, AdamW rearranges the
mechanism substantially---saturating drift, $B$-independent step noise, a
horizon-free threshold from weight decay---and leaves the exponent at $-1/2$. In
a linear-attention surrogate, parameter dimension leaves it there too. Endogenous
bias and horizon coupling both push the wrong way. What survives is total sample
exposure, for which we give an identity and a single-variable control, and the
plain possibility that three seeds were never going to resolve $-0.72$ from
$-0.5$.

Where that leaves us: the qualitative mechanism is probably right, the rate is
probably wrong though our replication falls short of establishing it, the
candidate explanations are narrowed and quantified rather than settled, and the
three measurements that would settle them---more seeds, teacher-forced scoring,
and the diversity control---are specified and runnable.

% ==================================================================
\bibliographystyle{plainnat}
\bibliography{references}

\appendix

% ==================================================================
\section{Proof of Theorem~\ref{thm:escape}}
\label{app:escape}

We prove (i); part (ii) is symmetric, with the inequalities reversed and $\mu_1$
in place of $\mu_0$. Let $\cF_t$ be generated by the first $t$ minibatches. On
the event that $u_s\in\cU$ for all $s\le t$, split the trajectory into drift and
martingale parts:
\[
u_T = u_0 + \underbrace{\sum_{t<T}\eta_t\E[h_t\mid\cF_t]}_{D_T}
+ \underbrace{\sum_{t<T}\eta_t(h_t-\E[h_t\mid\cF_t])}_{M_T}.
\]
In the zero-bias regime, Assumption~\ref{ass:main}(1)--(2) give
$\E[h_t\mid\cF_t]\ge\rho\mu_0$, so $D_T\ge\rho\mu_0S_T$. The martingale $M_T$ has
predictable quadratic variation bounded via \eqref{eq:batchvar} as
$\langle M\rangle_T = \sum_{t<T}\eta_t^2\Var(h_t\mid\cF_t)\le C_\eta V/B$, and by
Assumption~\ref{ass:main}(4) its increments are sub-exponential with parameter
$K_e$. Freedman's inequality then gives
$\Prob(M_T\le-\lambda)\le\exp(-\lambda^2/[2(\nu+K_e\lambda/3)])$ for any
$\nu\ge\langle M\rangle_T$.

Since $u_0=0$ and escape means $u_T\ge a$, it is enough that $D_T+M_T\ge a$.
Given $D_T\ge\rho\mu_0S_T$, that holds whenever
$M_T\ge a-\rho\mu_0S_T = -\lambda$ with $\lambda := \rho\mu_0S_T-a$, which is
positive exactly when $\rho>\rho^{+}_{\mathrm{esc}}$. Note that $\lambda$ here is
the actual slack between cumulative drift and escape radius, not a free tuning
parameter; applying Freedman with $\nu = C_\eta V/B$ gives the stated bound.
Since $\lambda = \mu_0S_T(\rho-\rho^{+}_{\mathrm{esc}})$, separating the two
branches of the denominator gives a failure probability at most
\[
\exp\Big(-c\min\Big\{\tfrac{B\mu_0^2S_T^2(\rho-\rho^{+}_{\mathrm{esc}})^2}{C_\eta V},\
\tfrac{\mu_0S_T(\rho-\rho^{+}_{\mathrm{esc}})}{K_e}\Big\}\Big)
\]
for a universal constant $c$. \qed

% ==================================================================
\section{Proof of Theorem~\ref{thm:matrixescape}}
\label{app:matrix}

\paragraph{Part (i): escaping needs one direction.}
Let $H_t := -\bar G_t$, so $U_{t+1} = U_t+\eta_tH_t$, and let
$u_t := \langle U_t-U_0,\Om^\star\rangle$. The projection is linear and
$\|\Om^\star\|_F = 1$, so $u_{t+1}-u_t = \eta_t\langle H_t,\Om^\star\rangle
=: \eta_th_t$ and $(u_t)$ obeys the scalar recursion. It remains to check
Assumption~\ref{ass:main} for $h_t$ with $v_+\mapsto v_\parallel$.

Drift: Assumption~\ref{ass:matrix}(1)--(2) in the zero-bias regime give
$\E[h_t\mid\cF_t] = \rho\mu(U_t)\ge\rho\mu_0$. Variance: for any unit $\Om$,
$\Var(\langle\hat G,\Om\rangle) =
\langle\Om,\mathrm{Cov}(\vecop\hat G)\Om\rangle\le\langle\Om,\Sigma\Om\rangle$,
so at $\Om^\star$ the per-example projected variance is at most $v_\parallel$.
Applying \eqref{eq:totalvar}--\eqref{eq:batchvar} to the projected per-example
gradients---legitimate because they are i.i.d.\ draws from the projected mixture
law by Definition~\ref{def:mixture}---gives $\Var(h_t\mid\cF_t)\le V_\parallel/B$.
Tails follow from Assumption~\ref{ass:matrix}(4) at $\Om^\star$. Finally, escape
from $\cU_a$ in the trace direction is implied by $u_T\ge a$, since
$\|U_T-U_0\|_F\ge\langle U_T-U_0,\Om^\star\rangle$ by Cauchy--Schwarz. So
Theorem~\ref{thm:escape}(i) applies with $V\mapsto V_\parallel$, with no union
bound over directions and hence no dependence on $D$.

\paragraph{Part (ii): staying put needs all of them.}
Confinement requires $\|U_t-U_0\|_F<a$ for all $t\le T$ simultaneously. Decompose
$U_T-U_0 = \mathbf D_T+\mathbf M_T$ with
$\mathbf D_T := \sum_{t<T}\eta_t\E[H_t\mid\cF_t]$ and
$\mathbf M_T := \sum_{t<T}\eta_t(H_t-\E[H_t\mid\cF_t])$, both in $\R^D$. By
Assumption~\ref{ass:matrix}(1), $\|\mathbf D_T\|_F\le\rho\mu_1S_T$.

To use a matrix martingale inequality on a vector-valued martingale, embed each
increment $\xi_t := \eta_t(H_t-\E[H_t\mid\cF_t])\in\R^D$ in the self-adjoint
dilation
\[
\mathcal D(\xi_t) := \begin{pmatrix}0 & \xi_t^\top\\ \xi_t & 0\end{pmatrix}
\in\R^{(D+1)\times(D+1)},
\]
which satisfies $\|\mathcal D(\xi)\|_{\mathrm{op}} = \|\xi\|_2$ and is linear in
$\xi$. So $\mathcal D(\mathbf M_T) = \sum_t\mathcal D(\xi_t)$ is a self-adjoint
matrix martingale with
$\|\mathcal D(\mathbf M_T)\|_{\mathrm{op}} = \|\mathbf M_T\|_2$, and its
predictable quadratic variation satisfies
\[
\Big\|\sum_{t<T}\E\big[\mathcal D(\xi_t)^2\mid\cF_t\big]\Big\|_{\mathrm{op}}
\;\le\; \sum_{t<T}\eta_t^2\,\Tr\,\mathrm{Cov}(H_t\mid\cF_t)
\;\le\; \frac{C_\eta V_\perp}{B},
\]
using $\mathcal D(\xi)^2 = \mathrm{diag}(\|\xi\|_2^2,\ \xi\xi^\top)$ and bounding
the operator norm of both blocks by the trace. With increments sub-exponential
with parameter $K_e$ in every direction by Assumption~\ref{ass:matrix}(4), the
matrix Freedman inequality of \citet{tropp2011freedman} gives
\[
\Prob\Big(\max_{t\le T}\|\mathbf M_t\|_2\ge\lambda'\Big)
\le (D+1)\exp\Big(-\frac{\lambda'^2}{2(C_\eta V_\perp/B+K_e\lambda'/3)}\Big),
\]
with the prefactor supplied by the ambient dimension of the dilation. On the
complement, $\|U_t-U_0\|_F\le\|\mathbf D_t\|+\|\mathbf M_t\|<a$ for all $t\le T$
with $\lambda' = a-\rho\mu_1S_T$. \qed

\paragraph{Where the dimension goes.}
The asymmetry is structural, and worth stating in one line: escaping needs
\emph{one} direction to move far, while staying needs \emph{every} direction to
stay close. That is why $v_\parallel$ appears in (i) and $\Tr\Sigma$ in (ii), and
why the undecided window widens by $\sqrt{D_{\mathrm{eff}}}$ in high dimension
while its endpoints keep the same batch-size scaling. We have not optimised the
$D+1$ factor; a chaining argument would likely replace it with
$\sqrt{\log D_{\mathrm{eff}}}$.

% ==================================================================
\section{Numbers, and one calibration}
\label{app:numbers}

For $q = 26$ and $L\sim\mathrm{Unif}\{4,\dots,10\}$, with $c$ the gradient scale
of \eqref{eq:gradform}:
\begin{center}
\begin{tabular}{lll}
\toprule
Quantity & Value & Source \\
\midrule
$\bar A^\star$ & $0.8926$ & Proposition~\ref{prop:bayesacc}, averaged over $L$ \\
$A_{\mathrm{chance}} = \E[1/L]$ & $0.1565$ & Corollary~\ref{cor:levels} \\
$\mu_0/c$ & $0.7361$ & Proposition~\ref{prop:mu0} \\
$v_+/c^2 = \E[\tfrac1L(1-\tfrac1L)]$ & $0.1296$ & Lemma~\ref{lem:variance}, $\rho\to0$ \\
$\sqrt{v_+}/\mu_0$ & $0.4892$ & scale-free by Remark~\ref{rem:scale} \\
\bottomrule
\end{tabular}
\end{center}

\paragraph{Calibration.}
Theorem~\ref{thm:scaling} gives $\rho_c(B)\approx0.489\,\kappa/\sqrt B$. Calibrate
at $B=256$, where the three-seed $\hat\rho_c = 0.175$, and you get
$\kappa = 0.175\sqrt{256}/0.489 = 5.72$. After that the theory predicts $\rho_c$
at every other $B$ with no freedom left:
\begin{center}
\begin{tabular}{lccccc}
\toprule
$B$ & 32 & 64 & 128 & 256 & 512 \\
\midrule
predicted $\rho_c$ ($\kappa = 5.72$) & 0.495 & 0.350 & 0.247 & 0.175 & 0.124 \\
measured $\hat\rho_c$ (3 seeds) & 0.759 & 0.465 & 0.305 & 0.175 & 0.099 \\
ratio & 1.53 & 1.33 & 1.23 & 1.00 & 0.80 \\
\bottomrule
\end{tabular}
\end{center}
The monotone trend in the last row is the exponent discrepancy of
Section~\ref{sec:exponent} wearing different clothes. A correct exponent would
give a flat row.

\paragraph{Reading the calibration through the AdamW surrogate.}
Comparing Theorem~\ref{thm:scaling} with \eqref{eq:adamrhoc} at $e_B\approx0$
identifies $\kappa$ with the critical standardized SNR,
$\kappa = \theta/\sqrt{1-\theta^2}$, so $\theta = \kappa/\sqrt{1+\kappa^2} =
0.985$ and $\sqrt{1-\theta^2} = 0.172$. For our schedule the binding branch of
\eqref{eq:thetadef} is the displacement one:
$S_T\approx600\cdot\tfrac12(3{\times}10^{-4}) + 5400\cdot\overline\eta_{\cos}
\approx0.09+0.84 = 0.93$, so $S_T^{-1} = 1.08\gg\lambda_\wdc = 0.01$ and
$\theta = a/S_T$. The calibrated constant then corresponds to an escape radius
$a = \theta S_T\approx0.92$ in logit units, which is about the scale on which a
softmax over $L\le10$ positions changes its argmax. Encouraging, and no more than
that: $\theta$ is close to $1$, where $\theta/\sqrt{1-\theta^2}$ is steep and the
inferred $a$ correspondingly sensitive, and the per-seed $\hat\rho_c$ at $B=256$
span $0.132$--$0.208$, which gives $a$ a range of roughly $0.89$--$0.94$.

\paragraph{Numerical floor and susceptibility.}
With $\varepsilon_{\mathrm A} = 10^{-8}$ and $\sqrt{v_+}/c = 0.360$, we get
$e_B = 6.3\times10^{-7}$ at $B=512$ and $1.6\times10^{-7}$ at $B=32$; with
$\theta = 0.985$, Theorem~\ref{thm:adamthreshold} predicts a slope of
$-0.5(1-3.7\times10^{-6})$. By Proposition~\ref{prop:susceptibility}, matching
$-0.72$ needs $\chi = 0.22$; over our sweep $N$ ranges from $1.92\times10^5$ to
$3.07\times10^6$, so the required improvement in the $\rho$-independent factor is
$16^{0.22} = 1.8\times$.

% ==================================================================
\section{Deferred material}
\label{app:deferred}

\subsection{Can the architecture even represent the circuit?}
\label{app:repr}

The scalar analysis presupposes that the architecture contains a
parameterization implementing the trace-copy family. We keep that question
separate from the optimization question and adopt both directions from the
literature rather than re-deriving them.

For sufficiency, standard constructions for two-layer attention-based copying
\citep{olsson2022incontext,ekbote2025what} apply. If the mechanism decomposes
into one attention operation moving a feature that encodes a matched trace
segment into the residual stream at the answer position, plus one readout turning
that into answer logits, with positional discrimination needed among at most
$2\ell_{\max}+C_0$ positions, then for $m\ge2$ layers and width
$d\ge C_1\log(2\ell_{\max}+C_0)$ the class contains such a parameterization with
answer error $O(1/T)$ on correctly-traced inputs.

For necessity, \citet{ekbote2025what} show a single attention layer cannot
represent an exact induction-head mechanism for general-order Markov structure
while two layers can, and our task needs the same two-hop composition: locate a
segment by content, then read an index out of it. We flag this as an adopted
result whose transfer we have not independently verified, since their setting
(arbitrary-order Markov chains) differs from word-indexing in ways that would
need checking. We treat it as an open technical gap.

Together these suggest $m\ge2$ suffices given adequate width, so our four-layer
model is unlikely to be facing a representability bottleneck. We stop short of
concluding that the transition must therefore be the scalar drift dynamics.
Ruling out one alternative does not establish another. The mechanism could be
drift on a trace-copy direction, the supervision-density effect of
Section~\ref{sec:competing}, a curriculum-like ordering effect from a reused
pool, or some combination.

\subsection{How many samples to see the sign of the drift}
\label{app:samplecomplexity}

A complementary way to feel the difficulty at small $\rho$ is to ask how many
i.i.d.\ gradient samples you need just to determine the \emph{sign} of the
population drift.

\begin{theorem}
\label{thm:samplecomplexity}
Let Assumption~\ref{ass:main} hold with net drift
$m(\rho) := \rho\mu_0-(1-\rho)\bar\epsilon>0$ at $u_0$. For the empirical mean of
$n$ i.i.d.\ per-example gradients to have the correct sign with probability
$1-\delta$, it suffices that
$n\ge C\max\{\sigma^2\log(1/\delta)/m(\rho)^2,\ K_e\log(1/\delta)/m(\rho)\}$, the
two branches being the Gaussian and sub-exponential regimes of Bernstein's
inequality under Assumption~\ref{ass:main}(4).
\end{theorem}

In the zero-bias regime this reads $n\gtrsim v_+\log(1/\delta)/(\rho^2\mu_0^2)$,
diverging as $\rho^{-2}$. That is the same qualitative picture as the
information-exponent literature \citep{benarous2021online,abbe2023sgd}, where a
weak population signal forces a diverging sample requirement. It is also the same
$\rho^{-2}$ that shows up in \eqref{eq:quadratic}: set $n = B$ and you recover
$\rho\gtrsim\sqrt{v_+}/(\mu_0\sqrt B)$. The two results are one inequality seen
from two sides.

\subsection{A reduced mean-field model}
\label{app:meanfield}

The scalar analysis follows one coordinate. The empirical transition plainly
involves several interacting quantities, so here is a reduced model with the
minimum number of order parameters that can produce a genuine bifurcation. This
is a modelling proposal, not a consequence of the transformer dynamics, which is
why it lives in an appendix.

Let $x_t\in[0,1]$ be trace-copy strength (the probability that the answer readout
attends to a trace segment rather than the prompt), $y_t\in[0,1]$ prompt-search
strength, and $z_t\ge0$ match precision, so that the probability of attending to
a segment whose character field equals $C$ is $\phi(z_t)$ with $\phi(0) = 1/L$ and
$\phi(\infty) = 1$.

The third parameter is not redundant, and this is the modelling point: a model
can attend to the trace ($x$ large) while matching badly ($z$ small), in which
case it copies an arbitrary index and performs at chance. Accuracy comes from the
product, not from either factor. We model
\begin{align*}
\dot x &= \alpha x(\rho\phi(z)\bar A^\star-\bar p)(1-x)-\alpha\lambda_xx,\\
\dot y &= \beta y(\bar A^\star-\bar p)(1-y)-\beta\lambda_yy,\\
\dot z &= \gamma\rho x(1-\phi(z))-\gamma\lambda_zz,
\end{align*}
with $\bar p$ the mean confidence in the chosen index. Match precision improves
only in proportion to $\rho x$---only when the model both attends to the trace
and the trace is correct---while $x$ grows only when $\phi(z)$ is already large.
That mutual dependence is how the system can support both a stable ``ignore the
trace'' state and a stable ``use the trace'' state.

Write $F_\rho$ for the right-hand side and $(x^0,y^0,z^0)$ for the trace-ignoring
fixed point with $x^0 = z^0 = 0$. Call the system non-degenerate at $\rho^\star$
if $D_{(x,z)}F_{\rho^\star}$ has a simple zero eigenvalue with right and left
eigenvectors $v,w$; $w^\top\partial_\rho F_{\rho^\star}\ne0$; and
$w^\top D^2F_{\rho^\star}[v,v]\ne0$. These are the standard transversality
conditions for a transcritical bifurcation, and we state them as checkable
primitives rather than assuming a normal form. The alternative---writing
$\dot x = a(\rho-\rho_c)x-bx^2$ and declaring a bifurcation---assumes the
conclusion.

Under those conditions the bifurcation is transcritical and
$x^\star\propto(\rho-\rho^\star)$. If the third condition fails while the cubic
term does not vanish, it is a pitchfork with exponent $1/2$. If the second fails,
the transition is imperfect and $x^\star$ is smooth. So measuring the exponent of
the normalised accuracy
$\widehat A(\rho) :=
(A(\rho)-A_{\mathrm{chance}})/(\bar A^\star-A_{\mathrm{chance}})$ against
$(\rho-\rho_c)$ discriminates the three cases. It needs denser sampling of $\rho$
near threshold than our grid has.

\subsection{Beyond word-indexing: a task family}
\label{app:family}

A fair worry is that the transition is a property of word-indexing rather than of
trace supervision. So we specify a family of four algorithmic tasks sharing the
surface form $\langle\text{prompt}\rangle\ \langle\text{trace}\rangle\
\texttt{:}\ \langle\text{answer}\rangle$, which lets the mixing, masking and
evaluation machinery stay identical across them.

\begin{center}
\small
\begin{tabular}{llcc}
\toprule
Task & Correct trace & Chance & Ceiling \\
\midrule
\texttt{word\_index} & \texttt{0m 1y 2n \dots} & 0.157 & 0.893 \\
\texttt{count\_char} & running tally \texttt{b0 a1 a2 c2 \dots} & 0.231 & 1.000 \\
\texttt{sort\_letters} & selection sort \texttt{ifkhb>b ifkh>f \dots} & 0.000 & 1.000 \\
\texttt{multiply} & \texttt{4*49=196 60*49=2940 196+2940=3136} & 0.002 & 1.000 \\
\bottomrule
\end{tabular}
\end{center}

Three design points matter here. First, the chance floor is the best constant
predictor, not a uniform guess. For word-index the two coincide, since
$\Prob(I^\star = 0) = \E[1/L]$, but for the others they do not, and the constant
predictor is the right choice because it is what a model ignoring its input
achieves. Getting this wrong inflates the apparent size of a transition. Second,
only word-index has a ceiling below one, which makes it the informative case for
reference-level comparisons and makes the others informative about the shape of
the transition. Third, since the floors differ by two orders of magnitude,
cross-task comparison has to use $\widehat A(\rho)$, which puts every task on a
common $[0,1]$ scale and leaves $\rho_c$ as the only thing left to read off.

% ==================================================================
\section{Reproducibility}
\label{app:repro}

Every sweep reads its defaults from a single configuration module and differs
only in the axis it varies on top of the correct-trace ratio: the ratio sweep,
the task, the batch size at fixed steps, and the batch size at fixed total sample
budget. Dataset pools are generated once and cached, keyed by task, supervision
condition, seed, context length and a fingerprint of the sampler's output, so
editing a task definition invalidates its cache instead of silently reusing data
from an earlier definition.

Two correctness guards are worth recording, because both caught real errors
during development. Examples exceeding the context length are counted and
reported rather than silently truncated---a truncated example loses its answer
token entirely and would be scored as a failure with no indication why. And the
loss mask produced by the cached representation is checked for exact equality
against an independently constructed reference mask, under all three supervision
conditions of Definition~\ref{def:conditions}.

No equivalent guard existed on the generation budget at evaluation time, which is
how the $35$-token cap of Section~\ref{sec:anomaly} went unnoticed for as long as
it did. The corresponding guard is an assertion that the gold continuation fits
inside the budget for every held-out example.

\end{document}
